%% 9-2 成分の取り出し方：直交性（第15回：組 V9-2。ページ制限なし。偶数ページ）
%% 「名前と操作」の型（tools/kata9-int.md）の柱の章。頭文字 kixc。場合分けの表9-2（A〜I）と流れ図。
%%   9-1 は I・E・F・B・C を予告して使い、9-3 は J（複素係数 c_n）の行を足して再掲する。
%% ほかの巻との約束：4-12:380「直交…第9巻では sin や cos の波どうしの直交…内積が0」→ ③④。
%%   予告カード 9-2「対称性で0になる積分（8-6 の場合G）と、内積（4-12）」→ ①⑧。8-6 では奇関数の名前なし→⑧で定義。
%%   組どうしの約束：8-13 の棒（L=0.1 m）の直線 200(−1)^{n+1}/(nπ)・三角形 800/(n²π²)sin(nπ/2)、弦（まん中で偶数番なし、1/5 で5番めなし）、
%%   方形波のひずみ率 48.3%、三角波 12.1%、基本波の実効値 0.900（使わず）、半波整流の平均 1/π・全波 2/π、a_2=−2/(3π)。
%%   数はすべて Python で検算した（scratchpad/v9-2/py/c4.py ほか）。
\makeatletter
%% ===== 「名前と操作」の型（tools/kata9-int.md、組 V9-2）。頭文字 kixc は章ごとに置きかえる =====
\colorlet{kixcNc}{bkOrange!85!black}% 名前の枠の線（点線）
\colorlet{kixcSc}{bkTheme}% 実際の操作の枠の線（実線）
\colorlet{kixcSb}{bkTheme!8}% 実際の操作の地
\tikzset{
  kixc dan/.style={font=\sffamily\bfseries\footnotesize, rounded corners=1.2mm, inner xsep=1.6mm, inner ysep=0.5mm, line width=0.9pt},
  kixc dan1/.style={kixc dan, draw=kixcNc, dash pattern=on 2.2pt off 1.4pt, fill=white, text=bkInk},
  kixc dan2/.style={kixc dan, draw=kixcSc, fill=kixcSb, text=bkInk},
  kixc dan3/.style={kixc dan, draw=bkTheme, fill=bkTheme, text=bkOn},
  kixc name/.style={draw=kixcNc, line width=0.9pt, dash pattern=on 2.2pt off 1.4pt, rounded corners=1.2mm, inner xsep=1.2mm, inner ysep=1.2mm, fill=white},
  kixc seki/.style={draw=kixcSc, line width=0.9pt, rounded corners=1.2mm, inner xsep=1.2mm, inner ysep=1.2mm, fill=kixcSb},
  kixc tag/.style={font=\sffamily\bfseries\footnotesize, text=bkInk, inner sep=0.4mm},
  kixc q/.style={diamond, aspect=5, draw=bkYellow!75!black, line width=0.9pt, fill=bkYellowBg, font=\sffamily\footnotesize, align=flush center, inner xsep=0.8mm, inner ysep=0.6mm, text=bkInk},
  kixc do/.style={rectangle, draw=\bkTT, line width=0.9pt, fill=white, rounded corners=0.6mm, font=\sffamily\footnotesize, align=flush center, inner xsep=1.5mm, inner ysep=1.2mm, text=bkInk},
  kixc go/.style={-{Stealth[length=2.6mm, width=2.2mm]}, line width=0.9pt, bkInk!75, shorten <=0.5mm, shorten >=0.5mm},
  kixc yn/.style={font=\sffamily\bfseries\footnotesize, text=bkInk, inner sep=0.5mm},
}
% 実際の操作の4種類の名前（実線の札の字）  1 ふつうの定積分　2 表を引く　3 部分分数　4 行列の計算
\def\kixc@op#1{\ifcase#1\or ふつうの定積分\or 表を引く\or 部分分数\or 行列の計算\fi}
% 段の札（例題の中で）  \kixcdan{1|2|3}
\newcommand\kixcdan[1]{\tikz[baseline=(d.base)]\node[kixc dan#1] (d) {\vphantom{あ}\ifcase#1\or 名前\or 直す\or 楽をする\fi};}
% 直す段の札に、どの道具かをそえる（例題の中で）  \kixcnaosu{1|2|3|4} →「直す：表を引く」
\newcommand\kixcnaosu[1]{\tikz[baseline=(d.base)]\node[kixc dan2] (d) {\vphantom{あ}直す：\kixc@op{#1}};}
% 実線の札だけ（表の「手順」の頭・本文）  \kixcop{1|2|3|4}
\newcommand\kixcop[1]{\tikz[baseline=(d.base)]\node[kixc dan2] (d) {\vphantom{あ}\kixc@op{#1}};}
% 段の見出し（本文で）  \kixcdanhead{1|2|3}{見出し}
\newcommand\kixcdanhead[2]{\par\addvspace{8pt}\Needspace{7\baselineskip}\noindent
  \tikz[baseline=(d.base)]\node[kixc dan#1, font=\sffamily\bfseries\small, inner ysep=0.8mm] (d) {\vphantom{あ}\ifcase#1\or ①\,名前の段\or ②\,直す段\or ③\,楽をする段\fi};%
  \hspace{0.7\zw}{\sffamily\bfseries\color{\bkTT}#2}\par\nobreak\vspace{3pt}\noindent\ignorespaces}
% 場合の札  \kixccase{C}（′ は \kixccase{A$'$}）
\newcommand\kixccase[1]{\tikz[baseline=(c.base)]\node[fill=bkTheme, text=bkOn, font=\sffamily\bfseries\small, rounded corners=0.8mm, inner xsep=1mm, inner ysep=0.6mm, minimum width=5.4mm] (c) {\vphantom{Ag}#1};}
% 例題の1行め  \kixcbaai{ことば}{9-2}{C}（場合が2つなら {D,B}：通った順）
\newcommand\kixcbaai[3]{{\sffamily\bfseries 場合：}#1（表#2の\kixc@cases#3,\@nil）}
\def\kixc@cases#1,#2\@nil{\kixccase{#1}\ifx\relax#2\relax\else と\kixc@cases#2\@nil\fi}
% 式の中の「名前」（点線）と「実際の操作」（実線）
\newcommand\kixcN[1]{\tikz[baseline=(n.base)]\node[kixc name] (n) {$\displaystyle #1$};}
\newcommand\kixcNt[1]{\tikz[baseline=(n.base)]{\node[kixc name] (n) {$\displaystyle #1$};\node[kixc tag, anchor=north] at ([yshift=-0.8mm]n.south) {名前};}}
\newcommand\kixcS[1]{\tikz[baseline=(n.base)]\node[kixc seki] (n) {$\displaystyle #1$};}
% 札つきの実線の枠  \kixcSt{1|2|3|4}{式}（札の字：ふつうの定積分・表を引く・部分分数・行列の計算）
\newcommand\kixcSt[2]{\tikz[baseline=(n.base)]{\node[kixc seki] (n) {$\displaystyle #2$};\node[kixc tag, anchor=north] at ([yshift=-0.8mm]n.south) {\kixc@op{#1}};}}
% 場合分けの表の題の帯  \kixccasetitle{9-2}{題}
\newcommand\kixccasetitle[2]{\par\addvspace{6pt}\Needspace{8\baselineskip}\noindent\begin{tikzpicture}[baseline=(t.base)]
  \node[fill=bkTheme, text=bkOn, font=\sffamily\bfseries\small, rounded corners=1mm, inner xsep=2mm, inner ysep=1mm] (t) {場合分けの表\,#1};
  \node[anchor=west, font=\sffamily\bfseries, text=\bkTT, inner sep=0pt] at ([xshift=2mm]t.east) {#2};
  \end{tikzpicture}\par\nobreak\vspace{3pt}}
% 名前と操作の対応表の小さな絵  \kixcpic{1〜7}
%  1 フーリエ係数（1周期で cos をかける）　2 複素係数（回る矢印）　3 フーリエ変換（−∞〜∞）
%  4 ラプラス変換（0から、e^{-st} で小さく）　5 逆ラプラス変換（s の平面の極）　6 伝達関数（箱）　7 e^{At}（行列とうずまき）
\newcommand\kixcpic[1]{\begin{tikzpicture}[x=1mm, y=1mm, baseline=(current bounding box.center)]
  \ifcase#1\or
  \draw[bkGray, line width=0.4pt] (0,4) -- (14,4);
  \draw[bkInk, line width=0.8pt] (1,4) -- (1,7) -- (7,7) -- (7,1) -- (13,1) -- (13,4);
  \draw[bkTheme, line width=0.8pt, dash pattern=on 1.6pt off 0.9pt, domain=1:13, samples=40, smooth] plot (\x,{4+2.6*cos(30*(\x-1))});
  \draw[bkInk, line width=0.5pt] (1,-0.6) -- (1,-1.6) -- (13,-1.6) -- (13,-0.6);
  \or
  \draw[bkGray, line width=0.4pt] (2,4) -- (12,4) (7,-0.5) -- (7,8.5);
  \draw[bkInk, line width=0.7pt] (7,4) circle (3.6);
  \draw[bkTheme, line width=1.1pt, -{Stealth[length=1.6mm]}] (7,4) -- ({7+3.6*cos(-40)},{4+3.6*sin(-40)});
  \draw[bkInk, line width=0.6pt, -{Stealth[length=1.3mm]}] ({7+5*cos(10)},{4+5*sin(10)}) arc (10:-60:5);
  \or
  \draw[bkInk, line width=0.6pt, {Stealth[length=1.4mm]}-{Stealth[length=1.4mm]}] (0,1) -- (14,1);
  \draw[bkInk, line width=0.8pt] (1.5,1) -- (5.5,1) -- (5.5,6.5) -- (8.5,6.5) -- (8.5,1);
  \draw[bkTheme, line width=0.8pt, dash pattern=on 1.6pt off 0.9pt, domain=1.5:12.5, samples=50, smooth] plot (\x,{1+2.2*(1+cos(80*(\x-7)))/2});
  \or
  \draw[bkInk, line width=0.6pt, -{Stealth[length=1.4mm]}] (1,1) -- (14,1);
  \draw[bkInk, line width=0.6pt] (1,0) -- (1,8.5);
  \fill[bkTheme!25, domain=1:13, samples=40, smooth] (1,1) -- plot (\x,{1+6.5*exp(-0.3*(\x-1))}) -- (13,1) -- cycle;
  \draw[bkTheme, line width=0.8pt, domain=1:13, samples=40, smooth] plot (\x,{1+6.5*exp(-0.3*(\x-1))});
  \or
  \draw[bkGray, line width=0.4pt] (0,4) -- (14,4) (10,-0.5) -- (10,8.5);
  \foreach \kixcx/\kixcy in {3/4,6.5/6.8,6.5/1.2} {\draw[bkRed, line width=1pt] ({\kixcx-1},{\kixcy-1}) -- ({\kixcx+1},{\kixcy+1}) ({\kixcx-1},{\kixcy+1}) -- ({\kixcx+1},{\kixcy-1});}
  \or
  \draw[bkInk, line width=0.6pt, -{Stealth[length=1.4mm]}] (0,4) -- (4,4);
  \draw[bkTheme, line width=0.8pt, fill=bkTheme!12] (4.4,1) rectangle (9.6,7);
  \node[font=\footnotesize, text=bkInk, inner sep=0pt] at (7,4) {$G$};
  \draw[bkInk, line width=0.6pt, -{Stealth[length=1.4mm]}] (10,4) -- (14,4);
  \or
  \draw[bkInk, line width=0.6pt] (1.6,0.5) -- (0.6,0.5) -- (0.6,7.5) -- (1.6,7.5) (5.4,0.5) -- (6.4,0.5) -- (6.4,7.5) -- (5.4,7.5);
  \foreach \kixcx/\kixcy in {2.3/2.3,4.7/2.3,2.3/5.7,4.7/5.7} {\fill[bkInk] (\kixcx,\kixcy) circle (0.5);}
  \draw[bkTheme, line width=0.8pt, domain=0:600, samples=80, smooth] plot ({10.5+3.4*exp(-0.0035*\x)*cos(\x)},{4+3.4*exp(-0.0035*\x)*sin(\x)});
  \fi
\end{tikzpicture}}
\makeatother
%% ===== この章の部品（頭文字 kixc） =====
\makeatletter
\newcommand\kixchead[1]{\par\addvspace{6pt}\Needspace{6\baselineskip}\noindent{\sffamily\bfseries\color{\bkTT}#1}\par\nobreak\vspace{2pt}\noindent\ignorespaces}
\newcommand\kixcbox[1]{\par\addvspace{3pt}\noindent\textbf{#1}\par\nobreak\vspace{2pt}\noindent\ignorespaces}
\newcommand\kixcframe[2][100mm]{\par\addvspace{4pt}{\centering\tikz\node[draw=\bkTT, line width=0.9pt, fill=bkTheme!6, rounded corners=1.5mm, inner xsep=6mm, inner ysep=2.5mm, align=center]{#2};\par}\addvspace{4pt}\noindent\ignorespaces}
\newcommand\kixca{\par\vspace{2pt}\noindent\rightskip=0pt plus 8em\relax\hangindent=2.6\zw\hangafter=0\relax}
\newcommand\kixcb{\par\vspace{2pt}\noindent\rightskip=0pt plus 8em\relax\hangindent=4.3\zw\hangafter=0\relax}
\newcommand\kixcs{\rightskip=0pt plus 8em\relax\hangindent=2.6\zw\hangafter=1\relax}
%% ===== 表計算・2つの世界の図の型（tools/kata9.md、組 V9-1） =====
\colorlet{kixcN}{bkOrange!16}% 1行め（成分の番号・周波数）のマス
\colorlet{kixcK}{bkGreen!16}% 2行め（成分の大きさ）のマス
\colorlet{kixcS}{bkYellow!45}% 合計・注目するマス
% 表計算の数式バー  \kixcfx{セル}{式}（式の中の $ は \$、^ は \^{}）
\newcommand\kixcfx[2]{\tikz[baseline=(f.base)]{\node[draw=bkLine, line width=0.6pt, fill=white, minimum height=5.5mm, inner xsep=2mm, font=\sffamily\footnotesize, minimum width=12mm] (c) {#1};
  \node[anchor=west, draw=bkLine, line width=0.6pt, fill=white, minimum height=5.5mm, inner xsep=2mm, minimum width=78mm, text width=74mm, align=left, font=\ttfamily\footnotesize] (f) at ([xshift=1mm]c.east) {#2};}}
\colorlet{kixcW}{bkBlue!80!black}% 時間の世界の線（波形・動き）
\colorlet{kixcF}{bkOrange!85!black}% 周波数の世界の棒
\colorlet{kixcS}{bkRed!80!black}% s の世界の点（×）
% グラフのスタイル（どれも scale only axis：はばと高さは軸の箱の大きさ）
\pgfplotsset{
  kixc base/.style={bkplot, scale only axis, width=47mm, height=28mm, axis lines=left, grid=none,
    every axis x label/.style={at={(axis description cs:0.5,0)}, anchor=north, yshift=-5mm, font=\sffamily\small, text=bkInk},
    every axis y label/.style={at={(axis description cs:0,1)}, anchor=south, yshift=1.5mm, font=\sffamily\small, text=bkInk}},
  kixc wave/.style={kixc base, extra y ticks={0}, extra y tick labels={0}, extra y tick style={grid=major, major grid style={bkInk!45, line width=0.5pt}}},
  kixc spec/.style={kixc base, ymin=0},
  kixc splane/.style={bkplot, scale only axis, width=47mm, height=47mm, axis equal image, grid=major,
    every axis x label/.style={at={(ticklabel* cs:1)}, anchor=north east, yshift=-0.5mm, font=\sffamily\small, text=bkInk},
    every axis y label/.style={at={(ticklabel* cs:1)}, anchor=south, yshift=1mm, font=\sffamily\small, text=bkInk}},
  kixc sig/.style={kixcW, line width=1.4pt, no markers, samples=241, smooth},
  kixc part/.style={bkGreen!45!black, line width=0.8pt, dash pattern=on 3pt off 1.5pt, no markers, samples=241, smooth},
  kixc bar/.style={ycomb, draw=kixcF, line width=3pt, no markers},
  kixc dot/.style={only marks, mark=*, mark size=1.3pt, kixcW}}
% 2つの世界の見出しと、行き来の矢印  \kixcworlds[見出しの高さ]{左の題}{右の題}{上の矢印（→）}{下の矢印（←）}
%     左の軸は at={(0mm,0mm)}、右の軸は at={(76mm,0mm)}（はば 47 mm）に置く。見出しは軸の上（ふつう 42 mm）
\tikzset{
  kixc world/.style={draw, line width=1pt, rounded corners=1.5mm, fill=white, font=\sffamily\bfseries\small, text=bkInk,
    minimum width=30mm, minimum height=7mm, inner xsep=2mm},
  kixc bridge/.style={-{Stealth[length=2.4mm]}, line width=1pt, bkInk!80},
  kixc blab/.style={font=\sffamily\footnotesize, text=bkInk, inner sep=0.8mm}}
\newcommand\kixcworlds[5][42mm]{%
  \node[kixc world, draw=kixcW] (kixcL) at (23.5mm,#1) {#2};
  \node[kixc world, draw=kixcF] (kixcR) at (99.5mm,#1) {#3};
  \draw[kixc bridge] ([yshift=1.3mm, xshift=1.5mm]kixcL.east) -- node[kixc blab, above] {#4} ([yshift=1.3mm, xshift=-1.5mm]kixcR.west);
  \draw[kixc bridge] ([yshift=-1.3mm, xshift=-1.5mm]kixcR.west) -- node[kixc blab, below] {#5} ([yshift=-1.3mm, xshift=1.5mm]kixcL.east);}
% s の世界の点（axis の中で）  \kixcpole[色]{実部}{虚部}（×）  \kixczero[色]{実部}{虚部}（○）
\newcommand\kixcpole[3][kixcS]{\draw[#1, line width=1.3pt] (axis cs:#2,#3) ++(-1.3mm,-1.3mm) -- ++(2.6mm,2.6mm) (axis cs:#2,#3) ++(-1.3mm,1.3mm) -- ++(2.6mm,-2.6mm);}
\newcommand\kixczero[3][kixcS]{\draw[#1, line width=1.1pt, fill=white] (axis cs:#2,#3) circle[radius=1.2mm];}
% 小さなグラフのわく（目もりなし）  \kixcmini{(x,y)}{はば}{高さ}
\newcommand\kixcaxes[3]{\draw[bkInk, line width=0.5pt, -{Stealth[length=1.6mm]}] #1 -- ++(#2,0);}
\makeatother

\bkchapter{9}{2}{成分の取り出し方：\\直交性}{\bklevel{5}}{成分の取り出し方：直交性}
\bkmichishirube{4-9,4-12,5-11,5-15,5-16,8-6,9-1}{9-3,9-4}

\begin{goal}
\begin{itemize}[label={\color{\bkTT}●}, leftmargin=1.6\zw, itemsep=0pt, topsep=1pt]
\item \bklevel{3}\enspace ちがう周波数の $\sin$ どうしをかけて1周期でたすと0になることを、表計算とグラフで確かめられる。
\item \bklevel{4}\enspace 直交性を使って、フーリエ係数が「かけて平均の2倍」で取り出せるわけを説明できる。
\item \bklevel{5}\enspace 場合分けの表9-2 と流れ図で、波の形に合った楽な求め方を選び、係数を計算できる。
\end{itemize}
\end{goal}

\begin{hatena}
\noindent\begin{minipage}[c]{0.46\linewidth}
\raggedright
ピアノでド・ミ・ソを同時にならすと、3つの音の波が重なって、図の波形のように、1つの波形になります。マイクがひろうのは、この重なった波だけです。\par\vspace{3pt}
\textbf{この波形から、ミの音の大きさだけを取り出すには、どんな計算をすればよいでしょう。}
\end{minipage}\hfill
\begin{minipage}[c]{0.52\linewidth}
\centering
\begin{tikzpicture}[x=1mm, y=1mm]
  \draw[bkGray, line width=0.4pt] (0,0) -- (62,0);
  \draw[bkBlue, line width=1pt, domain=0:60, samples=300, smooth, variable=\t] plot ({\t},{5*(sin(\t*24)+0.6*sin(\t*30)+0.3*sin(\t*36))});
  \draw[bkInk, line width=0.5pt, -{Stealth[length=1.6mm]}] (0,-11) -- (0,12);
  \draw[bkInk, line width=0.5pt, -{Stealth[length=1.6mm]}] (0,0) -- (64,0) node[right, font=\footnotesize, inner sep=1pt] {$t$};
  \draw[bkInk, line width=0.5pt] (60,-1) -- (60,1);
  \node[font=\footnotesize, anchor=north] at (60,-11) {0.02秒};
  \node[font=\footnotesize, anchor=south, align=center] at (30,12) {ド＋ミ＋ソの波形（1回のくり返し）};
\end{tikzpicture}
\end{minipage}
\end{hatena}

\noindent 9-1 で、フーリエ係数は「$\sin n\omega_0t$（または $\cos n\omega_0t$）をかけて、1周期で平均し、2倍する」と求まることを、表計算で見つけました。方形波に $\sin t$ をかけると正の値がたまり、$\sin2t$ をかけると0になりました。でも、なぜほかの成分は\textbf{ぴったり}0になるのでしょう。それがわかれば、はてな？のように、重なった波から、ほしい成分だけを取り出せます。

この章の主役は「\textbf{直交}（ちょっこう）」ということばです。4-12 で、2つの矢印の内積が0のとき「直交している」といいました。4-12 の最後には、「第9巻では、$\sin$ や $\cos$ の波どうしの『直交』を使って、複雑な波を単純な波に分けます。その計算の中心にあるのも、この『内積が0』という考えです」と書きました。この章で、その意味がわかります。

そして後半では、係数の求め方を、波の形ごとに「場合分けの表9-2」と流れ図にまとめます。係数は名前で、実際の計算のしかたは、波の形によって楽な道があるのです。

\begin{omoidasou}
\bkrecall{4-9}{積を和に直す式：$\sin\alpha\sin\beta=\frac12\bigl\{\cos(\alpha-\beta)-\cos(\alpha+\beta)\bigr\}$ など。}
\bkrecall{4-12}{内積 $\vec a\cdot\vec b=a_xb_x+a_yb_y$ が0なら、2つの矢印は直交。$\vec b$ の向きの成分は $\tfrac{\vec a\cdot\vec b}{|\vec b|^2}$ 倍。}
\bkrecall{5-15}{$\cos nt$・$\sin nt$（$n\geqq1$）を1周期で積分すると0。}
\bkrecall{8-6}{場合G：グラフを折り返したり半回転したりして重なるなら、積分は0になったり2倍になったりする。}
\bkrecall{9-1}{$b_n=\tfrac{2}{T}\int_0^Tf(t)\sin n\omega_0t\,\dd t$ は「かけて平均の2倍」という名前。方形波は $\frac{4}{n\pi}$（奇数番）。}
\end{omoidasou}

★★★★★の章ですが、目でわかる・手でわかるの段は、表計算と4-9 の公式で読めます。式でわかるの段は、例題を流れ図でたどりながら、少しずつ進めば大丈夫です。

\begin{yattemiyou}[ちがう波をかけて、たしてみよう]
表計算で、1周期（0° から360°）を30° ごとの12か所に分け、$\sin t$ と $\sin2t$ をかけてたします。
\begin{enumerate}[label={\bkcircled[bkGreen]{\arabic*}}, leftmargin=2.4\zw, itemsep=0pt, topsep=2pt]
\item A2〜A13 に0, 30, …, 330、B2 に \texttt{=SIN(RADIANS(A2))}、C2 に \texttt{=SIN(2*RADIANS(A2))}。
\item D2 に \texttt{=B2*C2}、E2 に \texttt{=B2*B2} を入れ、13行めまでコピーする。
\item 14行めに、D列とE列の合計（\texttt{=SUM(D2:D13)} など）を出す。
\end{enumerate}
\begin{bkcenter}
\bksheet{5}{19}{%
1 & \SetCell{c} $t$（度） & \SetCell{bg=kixcN} $\sin t$ & \SetCell{bg=kixcN} $\sin2t$ & \SetCell{c} $\sin t\sin2t$ & \SetCell{c} $\sin t\sin t$ \\
2 & 0 & 0 & 0 & 0 & 0 \\
3 & 30 & 0.5 & 0.866 & 0.433 & 0.25 \\
4 & 60 & 0.866 & 0.866 & 0.75 & 0.75 \\
5 & 90 & 1 & 0 & 0 & 1 \\
6 & 120 & 0.866 & $-0.866$ & $-0.75$ & 0.75 \\
7 & 150 & 0.5 & $-0.866$ & $-0.433$ & 0.25 \\
8 & 180 & 0 & 0 & 0 & 0 \\
9 & 210 & $-0.5$ & 0.866 & $-0.433$ & 0.25 \\
10 & 240 & $-0.866$ & 0.866 & $-0.75$ & 0.75 \\
11 & 270 & $-1$ & 0 & 0 & 1 \\
12 & 300 & $-0.866$ & $-0.866$ & 0.75 & 0.75 \\
13 & 330 & $-0.5$ & $-0.866$ & 0.433 & 0.25 \\
14 & 合計 & & & \SetCell{bg=kixcS} ア & \SetCell{bg=kixcS} イ \\}
\bktabcaption{D列が「ちがう波どうしの積」、E列が「同じ波どうしの積」。}
\end{bkcenter}
アとイをうめよう。D列とE列で、合計はどうちがうかな？
\end{yattemiyou}

\noindent D列は、正の値（0.433、0.75）と負の値（$-0.75$、$-0.433$）が、同じ数ずつあります。たすと、ア＝0 です。E列は2乗なので、どれも0以上です。
\begin{bkhenkei}[0.56]
0.25+0.75+1+0.75+0.25=3 & E列の0° から150° \\
3\times2=6 & イ：後半も同じ6つ \\
\end{bkhenkei}
合計にきざみ（$30°=\frac{\pi}{6}=0.5236$ ラジアン）をかけると、細い長方形の面積の合計（5-10）になり、定積分の近い値です。D列は $0\times0.5236=0$、E列は $6\times0.5236=3.1416$ で、これは $\pi$ です。
E列の値は、定積分 $\int_0^{2\pi}\sin^2t\,\dd t$ の近い値ですが、じつはこの場合はぴったり $\pi$ になります（②で確かめます）。

%% ---------------- 数学の道具 ----------------
\Needspace{14\baselineskip}
\bksection{doguu}{数学の道具}

\bkstage{me}{目でわかる}{ちがう波をかけると、打ち消す}{\bklevel{3}}

\bksubsection{1}{かけ算のグラフを見る}

\label{bk:9-2:kakeru}やってみようの $\sin t\times\sin2t$ と $\sin t\times\sin t$ を、なめらかなグラフにすると、図\ref{bk:9-2:fig-seki}のようになります。

\begin{bkfigure}
\begin{tikzpicture}[x=1mm, y=1mm]
  \foreach \kixcx/\kixcm/\kixct in {0/2/{$\sin t\times\sin2t$},68/1/{$\sin t\times\sin t$}} {%
    \begin{scope}[shift={(\kixcx,0)}]
    \node[font=\sffamily\footnotesize, anchor=south] at (28,13) {\kixct};
    \fill[bkBlue!25, domain=0:360, samples=150, variable=\t] (0,0) -- plot ({\t/6.5},{11*sin(\t)*sin(\kixcm*\t)}) -- (55.4,0) -- cycle;
    \draw[bkBlue, line width=1pt, domain=0:360, samples=150, variable=\t] plot ({\t/6.5},{11*sin(\t)*sin(\kixcm*\t)});
    \draw[bkInk, line width=0.5pt, -{Stealth[length=1.6mm]}] (0,0) -- (59,0) node[right, font=\footnotesize, inner sep=1pt] {$t$};
    \draw[bkInk, line width=0.5pt] (0,-9) -- (0,12);
    \foreach \kixcu/\kixcl in {13.85/{$\frac{\pi}{2}$},27.7/{$\pi$},55.4/{$2\pi$}} {\draw[bkInk, line width=0.5pt] (\kixcu,-0.8) -- (\kixcu,0.8); \node[font=\footnotesize, anchor=north] at (\kixcu,-9.5) {\kixcl};}
    \end{scope}}
  \node[font=\footnotesize, anchor=west, text=bkInk] at (-1,-20) {正と負の面積が同じ：合計0};
  \node[font=\footnotesize, anchor=west, text=bkInk] at (67,-20) {いつも0以上：平均 $\frac12$};
\end{tikzpicture}
\bkfigcaption{左：ちがう波どうしの積。0から $\pi$ の間でも、$\frac{\pi}{2}$ のまわりに半回転すると重なる形で、上と下の面積が打ち消す。右：同じ波どうしの積（$\sin^2t$）。0と1の間を2回ゆれ、平均は $\frac12$。}
\label{bk:9-2:fig-seki}
\end{bkfigure}

図\ref{bk:9-2:fig-seki}の左のグラフは、上と下に同じ形の山がならび、面積がちょうど打ち消し合います。たとえば0から $\pi$ の部分は、点 $(\frac{\pi}{2},\ 0)$ のまわりに半回転すると、もとのグラフにぴったり重なります。上の山が下の谷にうつるので、面積の合計は0です。これは 8-6 の場合G（「半回転すると重なるなら、積分は0」）と同じ考えです。

図\ref{bk:9-2:fig-seki}の右のグラフ $\sin^2t$ は、負になりません。$\sin^2t=\frac{1-\cos2t}{2}$（4-9 の2倍角の式）で、$\cos2t$ の部分は1周期で打ち消すので、平均は $\frac12$、1周期（長さ $2\pi$）の面積は $\frac12\times2\pi=\pi$ です。やってみようのE列の「合計×きざみ」3.1416 と同じです。

\kixchead{ほかの組み合わせも調べる}
同じように、いろいろな組み合わせで「1周期の積分」を数値で調べると、次のようになります。
\begin{bkcenter}
\begin{bktable}{colspec={Q[l,m,wd=42mm] Q[c,m,wd=20mm] Q[l,m,wd=42mm] Q[c,m,wd=20mm]}, row{1}={font=\sffamily\small}, rowsep=1.5pt, colsep=3pt}
かけたもの & 1周期の積分 & かけたもの & 1周期の積分 \\
$\sin t\times\sin2t$ & 0 & $\sin t\times\sin t$ & $\pi$ \\
$\sin2t\times\sin3t$ & 0 & $\sin3t\times\sin3t$ & $\pi$ \\
$\cos t\times\cos3t$ & 0 & $\cos2t\times\cos2t$ & $\pi$ \\
$\sin t\times\cos t$ & 0 & $\sin2t\times\cos2t$ & 0 \\
\end{bktable}
\bktabcaption{0から $2\pi$ までの積分。ちがう番号どうし、$\sin$ と $\cos$ どうしは0。同じ番号の $\sin$ どうし、$\cos$ どうしだけが $\pi$。}
\end{bkcenter}
「同じ番号の、同じ種類（$\sin$ と $\sin$、$\cos$ と $\cos$）」のときだけ0でなく、ほかはすべて0になります。この性質が、成分を取り出すかぎです。

%% ---------------- 手でわかる ----------------
\bkstage{te}{手でわかる}{式で確かめ、成分を取り出す}{\bklevel{4}}

\bksubsection{2}{積を和に直して、積分する}

\label{bk:9-2:sekiwa}表の結果を、式で確かめます。4-9 の積を和に直す式を使います（$m$、$n$ は1以上の整数）。
\begin{bkhenkei}[0.66]
\sin mt\sin nt=\frac12\bigl\{\cos(m-n)t-\cos(m+n)t\bigr\} & 4-9 の積和の式 \\
\int_0^{2\pi}\cos kt\,\dd t=\Bigl[\frac{\sin kt}{k}\Bigr]_0^{2\pi}=0 & $k$ が0でない整数なら（5-15） \\
\end{bkhenkei}
\kixchead{$m\neq n$ のとき}
$m-n$ も $m+n$ も0でない整数なので、2つの $\cos$ の積分はどちらも0です。
\[ \int_0^{2\pi}\sin mt\sin nt\,\dd t=\frac12(0-0)=0 \]
\kixchead{$m=n$ のとき}
$\cos(m-n)t=\cos0=1$ になるので、1つめの積分だけが残ります。
\begin{bkhenkei}[0.68]
\int_0^{2\pi}\sin^2nt\,\dd t=\frac12\Bigl(\int_0^{2\pi}1\,\dd t-\int_0^{2\pi}\cos2nt\,\dd t\Bigr) & 積和で $m=n$ \\
=\frac12(2\pi-0)=\pi & 1を積分すると長さ $2\pi$ \\
\end{bkhenkei}
$\cos$ どうし、$\sin$ と $\cos$ も、4-9 のほかの積和の式で同じように確かめられます（例題で使います）。まとめると、次のとおりです。
\kixcframe[118mm]{$\displaystyle\int_0^{2\pi}\sin mt\sin nt\,\dd t=\int_0^{2\pi}\cos mt\cos nt\,\dd t=\begin{cases}0 & (m\neq n)\\ \pi & (m=n)\end{cases}$\\[6pt]$\displaystyle\int_0^{2\pi}\sin mt\cos nt\,\dd t=0$\quad （どんな $m$、$n$ でも）}
ちがう番号の波どうしをかけて1周期で積分すると、\textbf{ぴったり0}になります。表計算で0になったのは、たまたまではありませんでした。

\bksubsection{3}{関数の内積と、直交}

\label{bk:9-2:naiseki}4-12 では、矢印 $\vec a=(a_1,\ a_2)$ と $\vec b=(b_1,\ b_2)$ の内積を $a_1b_1+a_2b_2$ と計算し、0なら直交とよびました。成分が3つなら $a_1b_1+a_2b_2+a_3b_3$ です。

やってみようの表のB列（$\sin t$ の12個の値）とC列（$\sin2t$ の12個の値）を、それぞれ「12個の成分をもつベクトル」と考えてみます。
\[ \vec a=(0,\ 0.5,\ 0.866,\ 1,\ \cdots,\ -0.5) \]
\[ \vec b=(0,\ 0.866,\ 0.866,\ 0,\ \cdots,\ -0.866) \]
D列は「成分どうしのかけ算」、その合計アは「内積 $\vec a\cdot\vec b$」そのものです。ア＝0 なので、この2つの12成分のベクトルは\textbf{直交}しています。

\begin{bkfigure}
\begin{tikzpicture}[x=1mm, y=1mm]
  % 左：矢印の内積
  \node[font=\sffamily\bfseries\footnotesize, text=\bkTT] at (22,30) {矢印（2成分）};
  \draw[bkInk, line width=0.5pt, -{Stealth[length=1.6mm]}] (0,0) -- (42,0);
  \draw[bkInk, line width=0.5pt, -{Stealth[length=1.6mm]}] (0,0) -- (0,26);
  \draw[bkBlue, line width=1.4pt, -{Stealth[length=2.4mm]}] (0,0) -- (30,12) node[right, font=\footnotesize, text=bkInk] {$\vec a$};
  \draw[bkOrange!85!black, line width=1.4pt, -{Stealth[length=2.4mm]}] (0,0) -- (-7,17.5) node[above, font=\footnotesize, text=bkInk] {$\vec b$};
  \node[font=\footnotesize, align=center] at (22,-8) {$a_1b_1+a_2b_2=0$\\（直角）};
  % まん中の矢印
  \draw[bkInk!75, line width=1pt, -{Stealth[length=2.4mm]}] (48,10) -- (60,10);
  \node[font=\sffamily\footnotesize, anchor=south, align=center] at (54,11) {成分を\\ふやす};
  % 右：関数の内積
  \node[font=\sffamily\bfseries\footnotesize, text=\bkTT] at (96,30) {関数（たくさんの成分）};
  \draw[bkGray, line width=0.4pt] (66,10) -- (126,10);
  \draw[bkBlue, line width=1pt, domain=0:360, samples=100, smooth, variable=\t] plot ({66+\t/6},{10+8*sin(\t)});
  \draw[bkOrange!85!black, line width=1pt, dash pattern=on 3pt off 1.5pt, domain=0:360, samples=100, smooth, variable=\t] plot ({66+\t/6},{10+8*sin(2*\t)});
  \foreach \kixck in {0,...,11} {\fill[bkInk] ({66+(\kixck*30)/6},{10+8*sin(\kixck*30)}) circle (0.6);}
  \node[font=\footnotesize, align=center] at (96,-8) {$\displaystyle\sum f_kg_k\Delta t\ \to\ \int_0^{2\pi}f\,g\,\dd t=0$};
\end{tikzpicture}
\bkfigcaption{左：2成分の矢印の内積。右：関数を「たくさんの点の値を並べたベクトル」と見ると、内積は「値どうしの積の合計」で、点を細かくした先は積の定積分になる。$\sin t$（青）と $\sin2t$（オレンジの点線）は直交している。}
\label{bk:9-2:fig-naiseki}
\end{bkfigure}

点を細かくしていくと、「値どうしの積の合計×きざみ」は、積の定積分 $\int f(t)g(t)\,\dd t$ に近づきます（5-10 の、細かく分けてたす考え）。そこで、1周期の関数 $f$ と $g$ について
\kixcframe[100mm]{$\displaystyle\int_0^{2\pi}f(t)\,g(t)\,\dd t$\quad を、$f$ と $g$ の\textbf{内積}とよぶ}
ことにします。内積が0のとき、2つの関数は\textbf{直交している}といいます。\bkidx{ちょっこう}{直交（関数の）}②の結果は、「ちがう番号の $\sin$・$\cos$ はすべて直交している」ということです。この性質を、三角関数の\textbf{直交性}（ちょっこうせい）といいます。

ベクトルの「長さの2乗」は自分との内積（$|\vec a|^2=\vec a\cdot\vec a$）でした。関数でも同じように考えると、$\sin nt$ の「長さの2乗」は $\int_0^{2\pi}\sin^2nt\,\dd t=\pi$ です。

\bksubsection{4}{成分を取り出す}

\label{bk:9-2:toridasu}4-12 では、矢印 $\vec a$ を、たがいに直交する2つの矢印 $\vec e_1$、$\vec e_2$ の和 $\vec a=c_1\vec e_1+c_2\vec e_2$ に分けるとき、係数は
\[ c_1=\frac{\vec a\cdot\vec e_1}{|\vec e_1|^2} \]
で求まりました。両辺を $\vec e_1$ と内積すると、$\vec e_2$ の項が直交で消えて、$\vec a\cdot\vec e_1=c_1|\vec e_1|^2$ だけが残るからです。フーリエ級数でも、まったく同じことをします。

\kixchead{$\sin$ だけでたした波の、$b_3$ を取り出す}
たとえば、$f(t)=b_1\sin t+b_2\sin2t+b_3\sin3t+b_4\sin4t+\cdots$ から、$b_3$ を取り出します。両辺に $\sin3t$ をかけて、0から $2\pi$ まで積分します。
\begin{bkhenkei}[0.66]
\int_0^{2\pi}f\sin3t\,\dd t & 両辺に $\sin3t$ をかけて積分 \\
=\textstyle b_1\int_0^{2\pi}\sin t\sin3t\,\dd t+b_2\int_0^{2\pi}\sin2t\sin3t\,\dd t & $b_1$・$b_2$ の項 \\
\quad\textstyle+b_3\int_0^{2\pi}\sin^23t\,\dd t+b_4\int_0^{2\pi}\sin4t\sin3t\,\dd t+\cdots & $b_3$・$b_4$…の項 \\
=b_1\times0+b_2\times0+b_3\times\pi+b_4\times0+\cdots & ②の直交性 \\
b_3=\frac{1}{\pi}\int_0^{2\pi}f(t)\sin3t\,\dd t & 両辺を $\pi$ でわる \\
\end{bkhenkei}
$b_3$ 以外の項は、直交性で0になって消えます。残った $b_3\times\pi$ の $\pi$ は、$\sin3t$ の「長さの2乗」です。だから、ほかの番号でも同じで
\[ b_n=\frac{\text{$f$ と $\sin nt$ の内積}}{\text{$\sin nt$ の長さの2乗}}=\frac{1}{\pi}\int_0^{2\pi}f(t)\sin nt\,\dd t \]
です。4-12 の $\frac{\vec a\cdot\vec e}{|\vec e|^2}$ と同じ形になりました。9-1 の「かけて平均の2倍」の正体はこれです。平均は $\frac{1}{2\pi}\int$ なので、$\frac{1}{\pi}\int$ は、ちょうど平均の2倍です（「長さの2乗」$\pi$ が、区間の長さ $2\pi$ の半分だから）。

\begin{bkfigure}
\begin{tikzpicture}[x=1mm, y=1mm]
  \node[kixc do, text width=22mm] (f) at (0,0) {重なった波 $f$};
  \node[kixc do, text width=24mm] (k) at (33,0) {$\sin3t$ をかける};
  \node[kixc do, text width=22mm] (s) at (66,0) {1周期で積分};
  \node[kixc do, text width=18mm] (b) at (96,0) {$\div\pi$ で $b_3$};
  \draw[kixc go] (f) -- (k);
  \draw[kixc go] (k) -- (s);
  \draw[kixc go] (s) -- (b);
  \node[font=\footnotesize, anchor=north, align=center] at (48,-7) {ほかの番号の成分は、ここで0になって消える（直交性）};
\end{tikzpicture}
\bkfigcaption{成分を取り出す手順。かける波を $\sin nt$ にかえれば、どの番号の成分でも取り出せる。}
\label{bk:9-2:fig-tejun}
\end{bkfigure}

\kixchead{はてな？の答え}
\bkhatenaref の波形を、計算しやすい数で考えます（本物のド・ミ・ソは約262・330・392\,Hz で、比がほぼ4：5：6 です。ここでは同じ比の200・250・300\,Hz とし、大きさも作った数にします）。
\[ f(t)=\sin(2\pi\times200t)+0.6\sin(2\pi\times250t)+0.3\sin(2\pi\times300t) \]
200・250・300\,Hz は、どれも50\,Hz の整数倍（4倍・5倍・6倍）なので、$f$ は $T=\frac{1}{50}=\qty{0.02}{s}$ でくり返します。ミ（250\,Hz）は5番めの成分です。ミの大きさを取り出すには、$\sin(2\pi\times250t)$ をかけて、1周期0.02秒で積分し、「長さの2乗」でわります。
\begin{bkhenkei}[0.62]
\int_0^{0.02}\sin^2(2\pi\times250t)\,\dd t=\frac{0.02}{2}=0.01 & 長さの2乗は周期の半分 \\
\int_0^{0.02}f\sin(2\pi\times250t)\,\dd t=0.6\times0.01 & ド・ソの項は直交性で0 \\
=0.006 & かけ算 \\
\frac{0.006}{0.01}=0.6 & 長さの2乗でわる \\
\end{bkhenkei}
ミの大きさ0.6が、ほかの音にじゃまされずに取り出せました。\textbf{ほしい周波数の $\sin$ をかけて、1周期で積分し、長さの2乗（周期の半分）でわる}。これが答えです。ドやソの大きさも、かける $\sin$ の周波数を200\,Hz・300\,Hz にかえれば、同じように1と0.3が取り出せます。

\kixcdanhead{1}{係数は、操作の名前}
ここで、名前と実際の操作を整理します。$b_n$ は「$f$ と $\sin n\omega_0t$ の内積を、長さの2乗でわる」という\textbf{操作の名前}です。記号を見ただけでは数は出ません。
\[ \kixcNt{b_n}=\kixcSt{1}{\frac{2}{T}\int_0^Tf(t)\sin n\omega_0t\,\dd t} \]
名前の段では、記号を読みます：\textbf{何に}（どの波 $f$）、\textbf{何をかけて}（$\sin n\omega_0t$ か $\cos n\omega_0t$）、\textbf{どこからどこまで}（1周期）、\textbf{何倍するか}（$\frac{2}{T}$）。

\kixcdanhead{2}{ふつうの定積分に直す}
実際に数を出すには、$f$ の式を入れて、ふつうの定積分を計算します。9-1 の方形波では、1周期を2つの区間に分けて計算しました。どんな波でも、まずは「区間ごとに $f$ の式を書き、それぞれ $\sin n\omega_0t$ をかけて積分してたす」と考えれば、計算できます。

\bkkokomade[ちがう番号の $\sin$・$\cos$ をかけて1周期で積分すると0になること（直交性）、それを使うと「かけて積分し、長さの2乗でわる」で成分が取り出せること。ここまでわかれば十分です。この先は、波の形ごとの楽な求め方（場合分けの表9-2）です。]

%% ---------------- 式でわかる ----------------
\Needspace{9\baselineskip}\bkstage{shiki}{式でわかる}{波の形ごとに、楽な求め方を選ぶ}{\bklevel{5}}

\bksubsection{5}{周期 $T$ の公式と、$\frac{2}{T}$・$\frac{a_0}{2}$ のわけ}

\label{bk:9-2:koshiki}周期が $T$ の波では、$\sin n\omega_0t$・$\cos n\omega_0t$（$\omega_0=\frac{2\pi}{T}$）が、0から $T$ までの1周期で直交します。②と同じ計算で、「長さの2乗」は周期の半分 $\frac{T}{2}$ になります（$\sin^2$ の平均が $\frac12$ だから）。だから、係数は「内積÷$\frac{T}{2}$」＝「内積×$\frac{2}{T}$」です。
\kixcframe[110mm]{$\displaystyle a_n=\frac{2}{T}\int_0^Tf(t)\cos n\omega_0t\,\dd t$\quad（$n\geqq0$）\\[6pt]$\displaystyle b_n=\frac{2}{T}\int_0^Tf(t)\sin n\omega_0t\,\dd t$\quad（$n\geqq1$）\\[6pt]$\displaystyle f(t)=\frac{a_0}{2}+\sum_{n=1}^{\infty}\bigl(a_n\cos n\omega_0t+b_n\sin n\omega_0t\bigr)$}

\kixchead{$\frac{a_0}{2}$ の $\frac12$ を、直交性で見る}
定数の項は、「$\cos0t=1$ の成分」です。1の「長さの2乗」は $\int_0^T1^2\,\dd t=T$ で、ほかの成分の $\frac{T}{2}$ の2倍です。だから、定数の項の正しい係数は「内積÷$T$」です。
\begin{bkhenkei}[0.6]
\text{定数の項}=\frac{1}{T}\int_0^Tf(t)\times1\,\dd t & 内積÷長さの2乗 $T$ \\
=\frac12\times\frac{2}{T}\int_0^Tf(t)\cos0t\,\dd t=\frac{a_0}{2} & $a_n$ の式を $n=0$ で使った形の半分 \\
\end{bkhenkei}
$a_n$ の式（÷$\frac{T}{2}$）を $n=0$ にもそのまま使うと、本当の値の2倍になってしまいます。それで $\frac12$ をつけて $\frac{a_0}{2}$ と書くのです（9-1 の「平均の2倍が $a_0$」と同じことの、別の見方です）。

\kixchead{1周期なら、どこからでもよい}
$f$ も $\sin$・$\cos$ も周期 $T$ でくり返すので、積分は $-\frac{T}{2}$ から $\frac{T}{2}$ でも同じです（9-1⑥）。$T=2\pi$ なら $\frac{2}{T}=\frac{1}{\pi}$ で、$\frac{1}{\pi}\int_{-\pi}^{\pi}$ の形もよく使います。

\bksubsection{6}{場合分けの表9-2 と流れ図}

\label{bk:9-2:basho}係数 $a_n$・$b_n$ は名前です。実際の計算は「ふつうの定積分」ですが、波の形によっては、積分しなくてすんだり、半分の区間ですんだりします。そこで、「こんなときは、こう求める」を表にしました（記号 A〜I）。9-1 で使った B・C・E・F・I もここに入っています。

表の「手順」の頭の札は、どの道具を使うかを表します。\kixcop{1} は直す段（ふつうの定積分で計算する）、\kixcdan{3} は楽をする段（計算の前に、形を楽なものに直す・積分をしない）です。楽をする段の場合は、そのあと直す段の場合へ進むことがあります（流れ図の「下の問いへ」）。

\kixchead{流れ図の使い方}
\begin{enumerate}[label={\bkcircled[\bkTT]{\arabic*}}, leftmargin=2.4\zw, itemsep=1pt, topsep=2pt]
\item 名前の段：どの波 $f$ か、周期 $T$、どの係数がほしいか（全部か、平均だけか）を書き出す。
\item 流れ図の上の問いから順に、「はい」か「いいえ」で答える。「はい」なら右の四角へ。
\item 「下の問いへ」と書いた四角に着いたら、矢印にそって左の列の問いへもどり、続ける。
\item 計算が終わったら、右の線をたどって「確かめ」をする。
\end{enumerate}
1回で全部の場合を覚えなくても大丈夫です。例題を解くたびに、表と流れ図にもどってきましょう。

\clearpage
%% ---- 表（左ページ） ----
\kixccasetitle{9-2}{フーリエ係数：どんなときに、どう求める？}
\begin{bkcenter}
\begin{bktable}{colspec={Q[c,m,9mm] Q[l,m,38mm] X[l,m] Q[c,m,10mm]}, row{2-Z}={font=\small}, rowsep=1.2pt, colsep=3pt}
記号 & こんなとき（見分け方） & 手順 & 例題 \\
\kixccase{I} & 式がなく、はかった値だけ & \kixcdan{3}\enspace 積分しない。$\frac{2}{N}\sum f_k\sin n\omega_0t_k$（「値×$\sin$」の平均の2倍。$N$ はデータの数） & 9 \\
\kixccase{A} & $\sin$・$\cos$ の積や累乗の式（$\sin^3t$、$\cos^2t$ など） & \kixcdan{3}\enspace 積分しない。4-9 の積和の式・2倍角の式で和に直し、係数を読む & 1 \\
\kixccase{G} & 平均（直流分）$\frac{a_0}{2}$ だけほしい & \kixcdan{3}\enspace 1周期のグラフの面積÷$T$ & 7 \\
\kixccase{E} & 棒・弦のように、0から $L$ だけで決まる形 & \kixcdan{3}\enspace 原点で半回転した形に広げる。$b_n=\frac{2}{L}\int_0^Lf\sin\frac{n\pi x}{L}\dd x$。積分は B・C で & 4・5 \\
\kixccase{F} & 周期 $T$ が $2\pi$ でない（50\,Hz など） & \kixcdan{3}\enspace $\omega_0=\frac{2\pi}{T}$、$\frac{2}{T}\int_0^T$ の式で。積分は D・B・C で & 6 \\
\kixccase{D} & 偶関数（たて軸で折り返して重なる）か奇関数（原点で半回転して重なる） & \kixcdan{3}\enspace 偶なら $b_n=0$、奇なら $a_n=0$。残りは半周期の2倍 $\frac{4}{T}\int_0^{T/2}$ & 2・3・6 \\
\kixccase{B} & 区間ごとに一定（方形波・パルス） & \kixcop{1}\enspace 区間に分けて定積分。$\sin n\pi=0$、$\cos n\pi=(-1)^n$ & 2・6 \\
\kixccase{C} & 区間ごとに1次式（のこぎり波・三角波） & \kixcop{1}\enspace 部分積分（5-16）。$t$ を微分する側に & 3・4・5 \\
\kixccase{H} & 区間ごとに $\sin$・$\cos$（整流した波） & \kixcop{1}\enspace 積和の式で和に直してから定積分 & 8 \\
\end{bktable}
\end{bkcenter}
\vspace{4pt}
\noindent{\small 楽をする段の札の場合（I・A・G・E・F・D）は、計算の前に形を楽なものに直します。E・F・D のあとは、直す段（B・C・H）で積分します。確かめのパーセバルの等式は⑪、9-3 では複素数の係数 $c_n$ の行（J・J$'$）を足します。}

\clearpage
%% ---- 流れ図（右ページ） ----
\begin{bkfigure}
\begin{tikzpicture}[x=1mm, y=1mm]
  % Ver2：共通の決まり2章の組み方（問いは左の列、「はい」は右、「いいえ」は下。右はしで集める線は使わない）
  \coordinate (QX) at (34,0); \coordinate (BX) at (76,0); \coordinate (CX) at (62,0);
  \node[kixc do, text width=96mm] (s) at (62,0) {記号を読む（名前の段）：どの波 $f$・周期 $T$・どの係数がほしいか};
  \coordinate (c0) at (s.south -| QX);
  \node[kixc q, below=7mm of c0] (q1) {式がなく、はかった値だけ？};
  \node[kixc q, below=7mm of q1] (q2) {$\sin$・$\cos$ の積・累乗の式？};
  \node[kixc q, below=7mm of q2] (q3) {平均だけほしい？};
  \node[kixc q, below=7mm of q3] (q4) {0〜$L$ だけの形（棒・弦）？};
  \node[kixc q, below=7mm of q4] (q5) {周期 $T$ が $2\pi$ でない？};
  \node[kixc q, below=7mm of q5] (q6) {偶関数か奇関数？};
  \node[kixc q, below=7mm of q6] (q7) {区間ごとに一定？};
  \node[kixc q, below=7mm of q7] (q8) {区間ごとに1次式？};
  \node[kixc do, text width=44mm, below=7mm of q8] (rH) {区間ごとの $\sin$・$\cos$：\\積和の式で和に直す\enspace\kixccase{H}};
  \node[kixc do, text width=52mm, anchor=west] (rI) at (q1 -| BX) {表計算で「値×$\sin$」の\\平均の2倍\enspace\kixccase{I}};
  \node[kixc do, text width=52mm, anchor=west] (rA) at (q2 -| BX) {積和の式で和に直し、\\係数を読む\enspace\kixccase{A}};
  \node[kixc do, text width=52mm, anchor=west] (rG) at (q3 -| BX) {グラフの面積÷$T$\enspace\kixccase{G}};
  \node[kixc do, text width=52mm, anchor=west] (rE) at (q4 -| BX) {半回転して広げる。$b_n=\frac{2}{L}\int_0^L\cdots$\enspace\kixccase{E}\\→ 積分は「区間ごとに一定？」から};
  \node[kixc do, text width=52mm, anchor=west] (rF) at (q5 -| BX) {$\omega_0=\frac{2\pi}{T}$、$\frac{2}{T}\int_0^T$\enspace\kixccase{F}\\→ 次の問いへ};
  \node[kixc do, text width=52mm, anchor=west] (rD) at (q6 -| BX) {0の係数をとばし、\\半周期の2倍\enspace\kixccase{D}\\→ 次の問いへ};
  \node[kixc do, text width=52mm, anchor=west] (rB) at (q7 -| BX) {区間に分けて定積分\enspace\kixccase{B}};
  \node[kixc do, text width=52mm, anchor=west] (rC) at (q8 -| BX) {部分積分\enspace\kixccase{C}};
  \path (rH.south) ++(0,-4mm) coordinate (e0);
  \node[font=\sffamily\footnotesize, text=bkInk, text width=124mm, align=flush center, anchor=north] (end) at (e0 -| CX) {I・A・G・B・C・H の四角で計算したら、最後に\textbf{確かめ}：\\何項かたしたグラフが、もとの波に近づくか（9-1）。2乗の平均が合うか（パーセバル）};
  \draw[kixc go] (c0) -- (q1);
  \foreach \kq/\kr in {q1/rI,q2/rA,q3/rG,q4/rE,q5/rF,q6/rD,q7/rB,q8/rC} {\draw[kixc go] (\kq) -- node[kixc yn, above=0.6mm] {はい} (\kr);}
  \foreach \kq/\kr in {q1/q2,q2/q3,q3/q4,q4/q5,q5/q6,q6/q7,q7/q8,q8/rH} {\draw[kixc go] (\kq) -- node[kixc yn, right=1mm] {いいえ} (\kr);}
\end{tikzpicture}
\bkfigcaption{フーリエ係数の流れ図。ひし形は問い、四角はすること。上の6つ（I・A・G・E・F・D）が楽をする段、下の3つ（B・C・H）が直す段。F・D のあとは左の列の次の問いへ、E のあとは「区間ごとに一定？」の問いへ進む（四角の中の → の先）。}
\label{bk:9-2:fig-flow}
\end{bkfigure}

\kixchead{流れ図の使い方：9-1 の熱の棒}
9-1⑧の「はじめどこも100\,℃の棒」をたどります。
\begin{enumerate}[label={\bkcircled[\bkTT]{\arabic*}}, leftmargin=2.4\zw, itemsep=1pt, topsep=2pt]
\item 名前の段：$f$ は0から $L$ で100、ほしいのは $\sin\frac{n\pi x}{L}$ の係数 $b_n$。
\item 式がなく、はかった値だけ？\enspace いいえ。\enspace $\sin$・$\cos$ の積・累乗の式？\enspace いいえ（定数100）。\enspace 平均だけ？\enspace いいえ。
\item 0〜$L$ だけの形？\enspace はい → \kixccase{E}：半回転して広げる。\par\nobreak\noindent $b_n=\frac{2}{L}\int_0^L100\sin\frac{n\pi x}{L}\dd x$ を求めればよい。
\item 右の線をたどって「区間ごとに一定？」へ。\enspace はい（0から $L$ で100）→ \kixccase{B}。
\end{enumerate}
こうして9-1 の例題6・7の「E と B」に着きます。

\kixchead{どの場合かな？}
次の係数は、表9-2 のどの場合でしょう。流れ図を上からたどって、先に自分で考えてから答えを見ましょう。
\begin{enumerate}[label={\bkcircled[\bkTT]{\arabic*}}, leftmargin=2.4\zw, itemsep=3pt, topsep=2pt]
\item $\cos^2t$ のフーリエ係数\par\nobreak\noindent 答え：\kixccase{A}。$\cos^2t=\frac12+\frac12\cos2t$ と書き直せる。
\item $-\pi<t<\pi$ で $f=|t|$ をくり返す三角波の $a_n$\par\nobreak\noindent 答え：\kixccase{D}→\kixccase{C}。偶関数で、区間ごとに1次式。
\item オシロスコープで16か所はかった電圧の波形の $b_1$\par\nobreak\noindent 答え：\kixccase{I}。式がなく、はかった値だけ。
\item 長さ $L$ の弦の、はじめの形（三角形）の $b_n$\par\nobreak\noindent 答え：\kixccase{E}→\kixccase{C}。0から $L$ だけの形で、区間ごとに1次式。
\end{enumerate}

\bksubsection{7}{楽をする段①：積分しない（A・G・I）}

\label{bk:9-2:A}\kixcdanhead{3}{計算する前に、楽な道をさがす}
流れ図の上の6つの問いは、③の\textbf{楽をする段}です。まず、積分そのものをしなくてすむ場合を見ます。

\kixchead{場合 A：公式で書き直して、係数を読む}
$f$ が、はじめから $\sin$・$\cos$ の積や累乗の式で書いてあるなら、4-9 の公式で $\sin$・$\cos$ の\textbf{和}に直すだけで、フーリエ級数の形になります。フーリエ級数の書き方は1通りしかない（せいかくに言うと）ので、和に直した式の係数が、そのままフーリエ係数です。たとえば、
\begin{bkhenkei}[0.56]
\cos^2t=\frac{1+\cos2t}{2}=\frac12+\frac12\cos2t & 4-9 の2倍角の式 \\
\frac{a_0}{2}=\frac12,\quad a_2=\frac12,\quad\text{ほかは0} & 係数を読む \\
\end{bkhenkei}
積分 $\frac{1}{\pi}\int_0^{2\pi}\cos^2t\cos2t\,\dd t$ を計算しても $\frac12$ になりますが、書き直すほうがずっと楽です。

\kixchead{場合 G：平均だけなら、面積÷周期}
ほしいのが平均（直流分）$\frac{a_0}{2}$ だけなら、「1周期のグラフの面積÷$T$」です（9-1⑥）。面積がわかる形なら、積分の計算もいりません。たとえば、高さ1の方形波（$1$ と $-1$）は、上と下の面積が同じなので、平均は0です。

\kixchead{場合 I：式がなく、はかった値だけなら、表計算でたす}
式がなく、1周期の中で等間かくに $N$ か所ではかった値 $f_0$、$f_1$、…、$f_{N-1}$ しかないときは、積分のかわりに和を使います。1周期で $N$ 個の値の「値×$\sin n\omega_0t_k$」の平均の2倍です。
\kixcframe[100mm]{$\displaystyle b_n\approx\frac{2}{N}\sum_{k=0}^{N-1}f_k\sin n\omega_0t_k,\qquad a_n\approx\frac{2}{N}\sum_{k=0}^{N-1}f_k\cos n\omega_0t_k$}
9-1⑤の表計算は、この式で計算していました。じつは、等間かくの点でも、②の直交性は、点の数が十分あれば\textbf{ぴったり}なり立ちます（やってみようの12点で、ア＝0、イ＝6＝$\frac{12}{2}$ でした）。この「とびとびの値のフーリエ係数」は、9-4 の DFT（離散フーリエ変換）につながります。

\bksubsection{8}{楽をする段②：偶関数・奇関数（D）}

\label{bk:9-2:gusu}8-6 の場合G では、「左右対称で $x$ の奇関数なら0」のように、形の対称性を使って積分を楽にしました。ここで、その形に正式に名前をつけます。

\begin{bkfigure}
\begin{tikzpicture}[x=1mm, y=1mm]
  \foreach \kixcx/\kixcname/\kixcfn in {0/{偶関数：$\cos t$、$|t|$、$t^2$}/0,68/{奇関数：$\sin t$、$t$、$t^3$}/1} {%
    \begin{scope}[shift={(\kixcx,0)}]
    \node[font=\sffamily\footnotesize, anchor=south] at (28,15) {\kixcname};
    \draw[bkInk, line width=0.5pt, -{Stealth[length=1.6mm]}] (-2,0) -- (59,0) node[right, font=\footnotesize, inner sep=1pt] {$t$};
    \draw[bkInk, line width=0.5pt, -{Stealth[length=1.6mm]}] (28,-13) -- (28,14);
    \end{scope}}
  \draw[bkBlue, line width=1.1pt, domain=-180:180, samples=100, smooth, variable=\t] plot ({28+\t/6.5},{10*cos(\t)});
  \draw[bkOrange!85!black, line width=0.8pt, -{Stealth[length=1.6mm]}] (38,8) to[bend right=40] (18,8);
  \node[font=\footnotesize, anchor=south] at (28,9.5) {};
  \node[font=\footnotesize, align=center, anchor=north] at (28,-14) {たて軸で折り返すと重なる\\$f(-t)=f(t)$};
  \draw[bkBlue, line width=1.1pt, domain=-180:180, samples=100, smooth, variable=\t] plot ({96+\t/6.5},{10*sin(\t)});
  \draw[bkOrange!85!black, line width=0.8pt, -{Stealth[length=1.6mm]}] (110,8) arc (30:210:16);
  \fill[bkInk] (96,0) circle (0.8);
  \node[font=\footnotesize, align=center, anchor=north] at (96,-14) {原点のまわりに半回転すると重なる\\$f(-t)=-f(t)$};
\end{tikzpicture}
\bkfigcaption{偶関数と奇関数。オレンジの矢印は、重ねるための動かし方。}
\label{bk:9-2:fig-gukisu}
\end{bkfigure}

すべての $t$ で $f(-t)=f(t)$ となる関数を、\textbf{偶関数}（ぐうかんすう）といいます。また、$f(-t)=-f(t)$ となる関数を、\textbf{奇関数}（きかんすう）といいます。\bkidx{ぐうかんすう}{偶関数・奇関数}$t^2$、$t^4$ のように偶数乗は偶関数、$t$、$t^3$ のように奇数乗は奇関数なので、この名前がつきました。$\cos t$ は偶関数、$\sin t$ は奇関数です（4-8：$\cos(-t)=\cos t$、$\sin(-t)=-\sin t$）。

\kixchead{かけ算の決まりと、積分}
\begin{bkcenter}
\begin{bktable}{colspec={Q[c,m,wd=26mm] Q[c,m,wd=26mm] Q[c,m,wd=26mm] Q[l,m,wd=40mm]}, row{1}={font=\sffamily\small}, rowsep=1.5pt, colsep=3pt}
$f$ & $g$ & $f\times g$ & $-a$ から $a$ までの積分 \\
偶 & 偶 & 偶 & 0から $a$ の2倍 \\
奇 & 奇 & 偶 & 0から $a$ の2倍 \\
偶 & 奇 & 奇 & 0（左右で打ち消す） \\
\end{bktable}
\bktabcaption{符号のかけ算（＋×＋、－×－、＋×－）と同じ決まり。奇関数×奇関数は偶関数。}
\end{bkcenter}
たとえば、$f$ が偶関数なら、$f\times\sin nt$ は「偶×奇＝奇」なので、$-\pi$ から $\pi$ の積分は0、つまり $b_n=0$ です。$f\times\cos nt$ は「偶×偶＝偶」なので、0から $\pi$ の積分の2倍になります。
\kixcframe[120mm]{偶関数：$b_n=0$、\quad $\displaystyle a_n=\frac{4}{T}\int_0^{T/2}f\cos n\omega_0t\,\dd t$\\[6pt]奇関数：$a_n=0$、\quad $\displaystyle b_n=\frac{4}{T}\int_0^{T/2}f\sin n\omega_0t\,\dd t$}
$\frac{4}{T}$ は「$\frac{2}{T}$ の2倍」（半分の区間の2倍）です。計算しなくてよい係数がわかり、残りも半分の区間ですみます。

\kixchead{原点の置き方で、偶奇は変わる}
同じ方形波でも、どこを $t=0$ にするかで、奇関数にも偶関数にもなります（図\ref{bk:9-2:fig-genten}）。9-1 の方形波（0から $\pi$ で1）は奇関数で、$\sin$ だけでした。時間を $\frac{\pi}{2}$ ずらして、$-\frac{\pi}{2}<t<\frac{\pi}{2}$ で1にすると偶関数になり、$\cos$ だけになります（例題2）。0と1をくり返す信号（9-1 の例題2）は、偶関数でも奇関数でもありません（平均 $\frac12$ をひくと奇関数）。

\begin{bkfigure}
\begin{tikzpicture}[x=1mm, y=1mm]
  \foreach \kixcx/\kixct in {0/{奇関数（$\sin$ だけ）},68/{偶関数（$\cos$ だけ）}} {%
    \begin{scope}[shift={(\kixcx,0)}]
    \node[font=\sffamily\footnotesize, anchor=south] at (28,10) {\kixct};
    \draw[bkInk, line width=0.5pt, -{Stealth[length=1.6mm]}] (-2,0) -- (59,0) node[right, font=\footnotesize, inner sep=1pt] {$t$};
    \draw[bkInk, line width=0.5pt, -{Stealth[length=1.6mm]}] (28,-9) -- (28,10);
    \end{scope}}
  \draw[bkOrange!85!black, line width=1.1pt] (-2,-6) -- (28,-6) -- (28,6) -- (58,6);
  \draw[bkOrange!85!black, line width=1.1pt] (66,-6) -- (81,-6) -- (81,6) -- (111,6) -- (111,-6) -- (126,-6);
  \node[font=\footnotesize, anchor=north] at (58,-1) {$\pi$};
  \node[font=\footnotesize, anchor=south west] at (111.6,0.6) {$\frac{\pi}{2}$};% Ver2：線に近すぎたので、線のない側へ
\end{tikzpicture}
\bkfigcaption{同じ方形波（周期 $2\pi$、たて軸は $t=0$）。左は原点で半回転すると重なる奇関数、右は時間を $\frac{\pi}{2}$ ずらしたもので、たて軸で折り返すと重なる偶関数。}
\label{bk:9-2:fig-genten}
\end{bkfigure}

\bksubsection{9}{直す段：区間ごとの形（B・C・H）}

\label{bk:9-2:BCH}\kixcdanhead{2}{ふつうの定積分に直す：区間ごとに式を書く}
楽をする段で区間や係数をへらしたら、残りをふつうの定積分で計算します。$f$ の形で、使う積分の道具が決まります。

\kixchead{場合 B：区間ごとに一定}
$\cos n\omega_0t$・$\sin n\omega_0t$ そのものの定積分です（9-1⑤）。区間のはしで $\sin n\pi=0$、$\cos n\pi=(-1)^n$ を使います。$\frac{\pi}{2}$ でくぎられる区間では、$\sin\frac{n\pi}{2}$ の値（$n=1$、2、3、4で1、0、$-1$、0 をくり返す）も使います（例題2）。

\kixchead{場合 C：区間ごとに1次式}
$t\cos nt$・$t\sin nt$ の形になるので、部分積分（5-16）で、$t$ を微分する側にします（9-1⑦）。$-\pi<t<\pi$ で $f=|t|$ をくり返す\textbf{三角波}なら、偶関数（D）なので $b_n=0$ で、$a_n$ は0から $\pi$ の2倍です。
\begin{bkhenkei}[0.62]
a_n=\frac{2}{\pi}\int_0^{\pi}t\cos nt\,\dd t & D：半周期の2倍。0から $\pi$ では $|t|=t$ \\
\int_0^{\pi}t\cos nt\,\dd t=\Bigl[\frac{t\sin nt}{n}\Bigr]_0^{\pi}-\frac{1}{n}\int_0^{\pi}\sin nt\,\dd t & 部分積分 \\
=0-\frac{1}{n}\times\frac{1-\cos n\pi}{n}=\frac{(-1)^n-1}{n^2} & $\sin n\pi=0$、$\cos n\pi=(-1)^n$ \\
a_n=\frac{2\bigl((-1)^n-1\bigr)}{\pi n^2} & $\frac{2}{\pi}$ をかける \\
\end{bkhenkei}
$n$ が奇数なら $(-1)^n-1=-2$ で $a_n=-\frac{4}{\pi n^2}$、偶数なら0です。平均は面積から $\frac{a_0}{2}=\frac{\pi}{2}$（G）。
\kixcframe[118mm]{$|t|=\dfrac{\pi}{2}-\dfrac{4}{\pi}\Bigl(\cos t+\dfrac{\cos3t}{9}+\dfrac{\cos5t}{25}+\cdots\Bigr)$\qquad（$-\pi<t<\pi$）}
方形波の係数は $\frac{1}{n}$ で小さくなりましたが、三角波は $\frac{1}{n^2}$ で、ずっと速く小さくなります。三角波には「とび」がなく、角があるだけだからです。とびのある形ほど、高い番号の成分がたくさん必要になります。

\begin{bkfigure}
\begin{tikzpicture}
\begin{axis}[bkplot, width=124mm, height=46mm, axis x line=bottom, axis y line=left, xmin=-200, xmax=560, ymin=0, ymax=3.6, xtick={-180,0,180,360,540}, xticklabels={$-\pi$,0,$\pi$,$2\pi$,$3\pi$}, ytick={1.5708,3.1416}, yticklabels={$\frac{\pi}{2}$,$\pi$}, xlabel={$t$}]
  \addplot[bkOrange!85!black, line width=1.2pt] coordinates {(-180,3.1416) (0,0) (180,3.1416) (360,0) (540,3.1416)};
  \addplot[bkBlue, line width=1pt, domain=-180:540, samples=300] {pi/2-4/pi*(cos(x)+cos(3*x)/9)};
\end{axis}
\end{tikzpicture}
\bkfigcaption{三角波（オレンジ）と、フーリエ級数の2項（$\cos t$ と $\cos3t$）までの和（青）。2項だけでも、かなりよく重なる。}
\label{bk:9-2:fig-sankaku}
\end{bkfigure}

\kixchead{場合 H：区間ごとに $\sin$・$\cos$}
スマートフォンの充電器などの電源の中では、交流をダイオードに通して、片方の向きだけを通す「整流」をします。\textbf{半波整流}（はんぱせいりゅう）した波は、$0<t<\pi$ で $\sin t$、$\pi<t<2\pi$ で0 です。係数の積分は $\int_0^{\pi}\sin t\cos nt\,\dd t$ のような形で、4-9 の積和の式
\[ \sin t\cos nt=\frac12\bigl\{\sin(1+n)t+\sin(1-n)t\bigr\} \]
で和に直してから積分します（例題8）。

\bksubsection{10}{棒と弦：半区間の展開（E）}

\label{bk:9-2:E}9-1 で、0から $L$ だけで決まる温度を、外側へ原点で半回転した形に広げて、周期 $2L$ の方形波にしました。この方法を一般の形に使います。

\kixchead{$b_n=\frac{2}{L}\int_0^L$ になるわけ}
広げた形は、周期 $T=2L$ の奇関数です。だから D で $a_n=0$、$b_n$ は半周期（0から $L$）の積分の2倍です。$\omega_0=\frac{2\pi}{2L}=\frac{\pi}{L}$ なので、
\begin{bkhenkei}[0.56]
b_n=\frac{4}{T}\int_0^{T/2}f\sin n\omega_0x\,\dd x & D の式（奇関数） \\
=\frac{4}{2L}\int_0^Lf(x)\sin\frac{n\pi x}{L}\,\dd x & $T=2L$、$\omega_0=\frac{\pi}{L}$ を入れる \\
=\frac{2}{L}\int_0^Lf(x)\sin\frac{n\pi x}{L}\,\dd x & $\frac{4}{2L}=\frac{2}{L}$ \\
\end{bkhenkei}
です。これを\textbf{半区間の展開}（はんくかんのてんかい）といいます。\bkidx{はんくかんのてんかい}{半区間の展開}外側の「見えない部分」は計算には出てきません。$\sin\frac{n\pi x}{L}$ は、どれも両はしで0なので、両はしを0\,℃に保つ棒や、両はしを止めた弦に、ちょうど合う形です。

\kixchead{いろいろな初めの形}
9-1 の棒（$L=\qty{0.1}{m}$）で、はじめの温度が一様でない場合も、同じように係数が求まります（例題4・じぶんでとこう5）。
\begin{bkcenter}
\begin{bktable}{colspec={Q[l,m,wd=34mm] Q[c,m,wd=34mm] Q[c,m,wd=15mm] Q[c,m,wd=15mm] Q[c,m,wd=15mm]}, row{1}={font=\sffamily\small}, rowsep=1.5pt, colsep=3pt}
はじめの形 & $b_n$ & $b_1$ & $b_2$ & $b_3$ \\
一様100\,℃（9-1） & $\dfrac{400}{n\pi}$（奇数） & 127.3 & 0 & 42.44 \\
直線 $100\frac{x}{L}$（右はし100） & $\dfrac{200(-1)^{n+1}}{n\pi}$ & 63.66 & $-31.83$ & 21.22 \\
三角形（まん中100） & $\dfrac{800}{n^2\pi^2}\sin\dfrac{n\pi}{2}$ & 81.06 & 0 & $-9.006$ \\
\end{bktable}
\bktabcaption{はじめの温度の形と、半区間の展開の係数。とびのある形ほど、$b_n$ がゆっくり小さくなる。}
\end{bkcenter}
答えは、どれも9-1 と同じく $T=\sum b_ne^{-\alpha(n\pi/L)^2t}\sin\frac{n\pi x}{L}$ です。$b_n$ だけがちがいます。

\kixchead{弦をはじく場所と、倍音}
ギターの弦（長さ $L$）のまん中を高さ $h$ だけつまんではなすと、はじめの形は三角形で、$b_n=\frac{8h}{n^2\pi^2}\sin\frac{n\pi}{2}$（例題5）。$n$ が偶数なら $\sin\frac{n\pi}{2}=0$ で、\textbf{偶数番の倍音（220\,Hz・440\,Hz…）が出ません}。まん中は、偶数番のゆれでは動かない点（節）だからです。はじく場所を、はしから $\frac{L}{5}$ にすると、$b_n$ は $\sin\frac{n\pi}{5}$ に比例し、5番め・10番め…の倍音が消えます。ギターをはじく場所で音色が変わるのは、このためです。

\begin{bkfigure}
\begin{tikzpicture}[x=1mm, y=1mm]
  \foreach \kixcx/\kixct/\kixcp in {0/{まん中ではじく}/0.5,68/{はしから $\frac15$ ではじく}/0.2} {%
    \begin{scope}[shift={(\kixcx,0)}]
    \node[font=\sffamily\footnotesize, anchor=south] at (27,20) {\kixct};
    \fill[bkInk] (0,10) circle (0.7) (54,10) circle (0.7);
    \draw[bkBlue, line width=1pt] (0,10) -- ({54*\kixcp},16) -- (54,10);
    \draw[bkInk, line width=0.5pt, -{Stealth[length=1.6mm]}] (2,-12) -- (58,-12) node[right, font=\footnotesize, inner sep=1pt] {$n$};
    \end{scope}}
  \foreach \kixcn/\kixch in {1/0.8106,2/0,3/0.0901,4/0,5/0.0324,6/0,7/0.0165} {\fill[bkOrange!85] ({2+6.5*\kixcn-1.8},-12) rectangle ({2+6.5*\kixcn+1.8},{-12+18*\kixch/0.8106}); \node[font=\footnotesize, anchor=north] at ({2+6.5*\kixcn},-12.6) {\kixcn};}
  \foreach \kixcn/\kixch in {1/0.7444,2/0.3011,3/0.1338,4/0.0465,5/0,6/0.0207,7/0.0246} {\fill[bkOrange!85] ({70+6.5*\kixcn-1.8},-12) rectangle ({70+6.5*\kixcn+1.8},{-12+18*\kixch/0.8106}); \node[font=\footnotesize, anchor=north] at ({70+6.5*\kixcn},-12.6) {\kixcn};}
\end{tikzpicture}
\bkfigcaption{はじめの形と、係数の大きさ $|b_n|$ の棒グラフ（スペクトル）。まん中ではじくと偶数番が0、はしから $\frac15$ ではじくと5番めが0。}
\label{bk:9-2:fig-hajiku}
\end{bkfigure}

\bksubsection{11}{パーセバルの等式：電力は成分ごとの和}

\label{bk:9-2:parseval}直交性からは、もう1つ大事な式が出てきます。$f$ の2乗の平均を考えます。$f=\frac{a_0}{2}+\sum(a_n\cos n\omega_0t+b_n\sin n\omega_0t)$ を2乗すると、たくさんの「かけ算の項」が出ますが、1周期で平均すると、ちがう成分どうしの積は直交性で0になります。残るのは同じ成分の2乗だけで、$\cos^2$・$\sin^2$ の平均は $\frac12$ です。
\kixcframe[110mm]{$\displaystyle\frac{1}{T}\int_0^Tf(t)^2\,\dd t=\Bigl(\frac{a_0}{2}\Bigr)^2+\frac12\sum_{n=1}^{\infty}\bigl(a_n^2+b_n^2\bigr)$}
これを\textbf{パーセバルの等式}といいます。\bkidx{ぱーせばるのとうしき}{パーセバルの等式}電圧 $f$ を抵抗にかけると、電力は $f^2$ に比例します。だから、この式は「\textbf{全体の電力は、成分ごとの電力の和}」という意味です（4-12 の三平方の定理 $|\vec a|^2=a_1^2+a_2^2$ の、関数版です）。

\kixchead{方形波で確かめる}
高さ1の方形波は、2乗するといつも1なので、左辺は1です。右辺は $b_n=\frac{4}{n\pi}$（奇数）から
\begin{bkhenkei}[0.56]
1=\frac12\Bigl\{\Bigl(\frac{4}{\pi}\Bigr)^2+\Bigl(\frac{4}{3\pi}\Bigr)^2+\cdots\Bigr\} & パーセバルの等式 \\
1=\frac{8}{\pi^2}\Bigl(1+\frac19+\frac{1}{25}+\cdots\Bigr) & $\frac12\times\frac{16}{\pi^2}=\frac{8}{\pi^2}$ でくくる \\
1+\frac19+\frac{1}{25}+\frac{1}{49}+\cdots=\frac{\pi^2}{8}=1.2337 & 両辺に $\frac{\pi^2}{8}$ をかける \\
\end{bkhenkei}
奇数の2乗の逆数の和が $\frac{\pi^2}{8}$ になりました（Python で1000項たすと1.2335）。方形波の電力のうち、基本波がもつのは $\frac{8}{\pi^2}=0.811$、つまり約81\% で、残りの約19\% は高調波がもっています。

\kixchead{ひずみ率}
電気の世界では、波が $\sin$ からどれだけくずれているかを、\textbf{ひずみ率}（全高調波ひずみ率）で表します。高調波の電力の合計を基本波の電力でわり、平方根をとったものです。
\kixcframe[100mm]{ひずみ率$=\sqrt{\dfrac{\text{高調波の電力の合計}}{\text{基本波の電力}}}=\dfrac{\sqrt{b_2^2+b_3^2+\cdots}}{b_1}$\\[3pt]\small\sffamily （$a_n$ があるときは、$a_n^2$ も同じようにたす）}
方形波では $\sqrt{1.2337-1}=\sqrt{0.2337}=0.483$、つまり48.3\% です（例題10）。三角波は係数が $\frac{1}{n^2}$ で小さくなるので12.1\%、きれいな $\sin$ の波は0\% です。

\Needspace{14\baselineskip}
\begin{motto}[5]
\kixcbox{フーリエ係数は、いちばん近い近似}
フーリエ級数を $N$ 項でやめたものは、「$\sin$・$\cos$ を $N$ 個まで使ってできる形のうちで、もとの $f$ に（2乗の平均の意味で）いちばん近いもの」になっています。4-12 で、点から平面へいちばん近い点は「直交する方向に下ろした足（正射影）」でしたが、それと同じことが関数でも起きているのです。たす項をふやしても、それまでの係数は変わりません（直交しているので、新しい項が古い項にじゃまをしない）。この考えを\textbf{最小2乗近似}といい、データの近似（3-14）や、画像・音の圧縮（9-12）の土台です。
\end{motto}

\begin{motto}[5]
\kixcbox{$\sin$ は「2回微分」の固有ベクトル}

\[ \frac{\dd^2}{\dd x^2}\sin\frac{n\pi x}{L}=-\Bigl(\frac{n\pi}{L}\Bigr)^2\sin\frac{n\pi x}{L} \]
です。「2回微分する」という操作をしても、形が変わらず何倍かになるだけです。7-8 で、行列 $A$ をかけても向きが変わらず何倍かになるベクトルを固有ベクトルとよびました。$\sin\frac{n\pi x}{L}$ は、「2回微分」の固有ベクトル（固有関数）です。熱方程式が $\sin$ の形に分けると解けたのは、このためです。9-11 では、棒を9点に区切った「熱の1歩の行列」の固有ベクトルが、本当に $\sin$ の形になることを確かめます。
\end{motto}

\Needspace{12\baselineskip}
\begin{seikaku}
\kixcbox{フーリエ級数の書き方は1通り}
$\sin n\omega_0t$・$\cos n\omega_0t$（と定数）は、どの2つも直交していて、しかも「足りないものがない」（どんな形も、これらのたし算で、2乗の平均の意味でいくらでも近くできる）ことが知られています。この性質を\textbf{完全性}といいます（証明は専門書で）。直交していて完全なので、1つの $f$ を書き表す係数の組は1通りに決まります。場合 A で、公式で書き直した係数をそのままフーリエ係数とみてよいのは、このためです。専門書では、関数の内積を $\langle f,\ g\rangle=\int f\,g\,\dd t$ と書きます。
\end{seikaku}

\begin{kotoba}
\small
\bkword{関数の内積}{1周期で $\int f\,g\,\dd t$。値どうしの積の合計の、点を細かくした先。}
\bkword{直交・直交性}{内積が0のとき直交。ちがう番号の $\sin$・$\cos$ がどれも直交していること。}
\bkword{偶関数・奇関数}{$f(-t)=f(t)$ が偶関数（たて軸で折り返して重なる）、$f(-t)=-f(t)$ が奇関数（原点で半回転して重なる）。}
\bkword{半区間の展開}{0から $L$ だけの形を、原点で半回転して広げ、$\sin\frac{n\pi x}{L}$ だけで書くこと。}
\bkword{整流}{交流の、片方の向きだけを通したり、向きをそろえたりすること。}
\bkword{パーセバルの等式}{2乗の平均は、成分の2乗の平均の和。}
\bkword{ひずみ率}{高調波の電力の合計÷基本波の電力の平方根。}
\end{kotoba}

\bkkigocards{\bkkigocard{$(-1)^n$}{マイナス・イチの・エヌじょう}{$-1$ を $n$ 回かけた数。$n$ が偶数なら1、奇数なら $-1$。$\cos n\pi$ の値}{minus one to the $n$-th power}{9-2}}

\begin{taisetsu}
\kixcbox{直交性}
ちがう番号の $\sin n\omega_0t$・$\cos n\omega_0t$ は、かけて1周期で積分すると0（直交）。同じものどうしは $\frac{T}{2}$（長さの2乗）。
\kixcbox{係数は名前。内積÷長さの2乗}
\[ \kixcNt{b_n}=\kixcSt{1}{\frac{2}{T}\int_0^Tf(t)\sin n\omega_0t\,\dd t} \]
\kixcdan{1}\enspace 記号を読む：どの波 $f$・周期 $T$・どの係数か\par
\kixcdan{3}\enspace 計算の前に：はかった値だけ（I）？\ 公式で書ける（A）？\ 平均だけ（G）？\ 0〜$L$（E）？\ 周期（F）？\ 偶奇（D）？\par
\kixcnaosu{1}\enspace 区間ごとに一定（B）・1次式（C）・$\sin$ の形（H）
\kixcbox{場合分けの表9-2（まとめ）}
\begin{bktable}{colspec={Q[c,m,9mm] Q[l,m,40mm] X[l,m]}, row{2-Z}={font=\small}, rowsep=1pt, colsep=3pt}
記号 & こんなとき & 直し方 \\
\kixccase{I} & はかった値だけ & $\frac{2}{N}\sum f_k\sin n\omega_0t_k$ \\
\kixccase{A} & $\sin$・$\cos$ の積・累乗 & 積和の式で和に直し、係数を読む \\
\kixccase{G} & 平均だけ & 1周期の面積÷$T$ \\
\kixccase{E} & 0から $L$ だけの形 & $b_n=\frac{2}{L}\int_0^Lf\sin\frac{n\pi x}{L}\dd x$ \\
\kixccase{F} & 周期が $2\pi$ でない & $\omega_0=\frac{2\pi}{T}$、$\frac{2}{T}\int_0^T$ \\
\kixccase{D} & 偶関数・奇関数 & 偶なら $b_n=0$、奇なら $a_n=0$。$\frac{4}{T}\int_0^{T/2}$ \\
\kixccase{B} & 区間ごとに一定 & 区間に分けて定積分 \\
\kixccase{C} & 区間ごとに1次式 & 部分積分 \\
\kixccase{H} & 区間ごとに $\sin$・$\cos$ & 積和の式で和に直して定積分 \\
\end{bktable}
\kixcbox{パーセバルの等式}
2乗の平均＝$\bigl(\frac{a_0}{2}\bigr)^2+\frac12\sum(a_n^2+b_n^2)$。全体の電力は成分ごとの電力の和。
\end{taisetsu}
%% ---------------- 例題 ----------------
\bksection{reidai}{例題でたしかめよう}

\noindent 例題の「考え方」の1行めには、場合分けの表9-2 のどの場合かを、流れ図で通った順に書きました。例題1〜10を、図\ref{bk:9-2:fig-flow}の流れ図の上から順にたどると、次のようになります（書いていない問いは「いいえ」）。じぶんでとこうも同じ道すじです。
\begin{bkcenter}
\begin{bktable}{colspec={Q[c,m,10mm] Q[l,m,30mm] X[l,m] Q[c,m,18mm]}, row{2-Z}={font=\small}, rowsep=0.8pt, colsep=3pt}
例題 & 題 & 「はい」と答えた問い & 着いた所 \\
1 & $\sin^3t$ & $\sin$・$\cos$ の積・累乗の式？ & \kixccase{A} \\
2 & 偶関数の方形波 & 偶関数か奇関数？ → 区間ごとに一定？ & \kixccase{D}\kixccase{B} \\
3 & 三角波 & 偶関数か奇関数？ → 区間ごとに1次式？ & \kixccase{D}\kixccase{C} \\
4 & 直線の初めの温度 & 0〜$L$ だけの形？ → 区間ごとに1次式？ & \kixccase{E}\kixccase{C} \\
5 & まん中ではじいた弦 & 0〜$L$ だけの形？ → 区間ごとに1次式？ & \kixccase{E}\kixccase{C} \\
6 & 50\,Hz の方形波 & 周期 $T$ が $2\pi$ でない？ → 偶関数か奇関数？ → 区間ごとに一定？ & \kixccase{F}\kixccase{D}\kixccase{B} \\
7 & 整流した波の平均 & 平均だけほしい？ & \kixccase{G} \\
8 & 半波整流の $a_2$ & （すべていいえ） & \kixccase{H} \\
9 & 12個のデータ & 式がなく、はかった値だけ？ & \kixccase{I} \\
10 & ひずみ率 & 係数はわかっている（流れ図は使わない） & ― \\
\end{bktable}
\bktabcaption{例題の流れ図の道すじ。}
\end{bkcenter}

\begin{reidai}[公式で書き直す]
$f(t)=\sin^3t$ のフーリエ係数を求めましょう（周期 $2\pi$）。\par
\kangaekata \kixcbaai{$\sin$ の3乗の式}{9-2}{A}\par
\kixcdan{1}\enspace $b_n=\frac{1}{\pi}\int_0^{2\pi}\sin^3t\sin nt\,\dd t$ など。\par
\kixcdan{3}\enspace 積分しない。4-9 の公式で $\sin^3t$ を $\sin$ の和に書き直し、係数を読む。\par
\shiki\kixcs $\sin^3t=\sin t\times\sin^2t=\sin t\times\frac{1-\cos2t}{2}$\quad （2倍角の式）\par
\kixca $=\frac12\sin t-\frac12\sin t\cos2t$\par
\kixca $\sin t\cos2t=\frac12\bigl\{\sin3t+\sin(-t)\bigr\}=\frac12\sin3t-\frac12\sin t$\quad （積和の式）\par
\kixca $\sin^3t=\frac12\sin t-\frac14\sin3t+\frac14\sin t=\frac34\sin t-\frac14\sin3t$\par
\begin{center}\begin{tikzpicture}
\begin{axis}[bkplot, axis x line=bottom, width=104mm, height=36mm, xmin=0, xmax=380, ymin=-1.1, ymax=1.25, xtick={90,180,270,360}, xticklabels={$\frac{\pi}{2}$,$\pi$,$\frac{3\pi}{2}$,$2\pi$}, ytick={-1,0,1}, xlabel={$t$}, legend style={font=\sffamily\footnotesize, at={(0.5,1.04)}, anchor=south, draw=bkLine, inner sep=1pt}, legend columns=3, legend cell align=left]
  % Ver2：目もりの字を曲線の下（グラフの下のはし）へ、凡例を上へ移した（字と線・凡例と軸の名前が重なっていた）
  \draw[bkInk!70, line width=0.5pt] (axis cs:0,0) -- (axis cs:375,0);
  \addplot[bkBlue, line width=1.3pt, domain=0:360, samples=120] {sin(x)^3};
  \addplot[bkOrange!85!black, line width=0.9pt, dash pattern=on 3pt off 1.5pt, domain=0:360, samples=120] {0.75*sin(x)};
  \addplot[bkGreen!70!black, line width=0.9pt, dash pattern=on 1pt off 1pt, domain=0:360, samples=120] {-0.25*sin(3*x)};
  \legend{$\sin^3t$,$\frac34\sin t$,$-\frac14\sin3t$}
\end{axis}
\end{tikzpicture}\end{center}
\kixca 2つの点線をたすと、青い線（$\sin^3t$）になる\par
\kotae $b_1=\frac34$、$b_3=-\frac14$、ほかはすべて0
\end{reidai}

\begin{jibun}
$f(t)=\cos^3t$ のフーリエ係数を求めましょう（周期 $2\pi$）。
\bkjans{$\cos^3t=\cos t\times\frac{1+\cos2t}{2}=\frac12\cos t+\frac12\cos t\cos2t$。$\cos t\cos2t=\frac12(\cos3t+\cos t)$ なので、$\cos^3t=\frac34\cos t+\frac14\cos3t$。$a_1=\frac34$、$a_3=\frac14$、ほかは0。}
\end{jibun}

\begin{reidai}[偶関数の方形波]
$-\frac{\pi}{2}<t<\frac{\pi}{2}$ で1、$\frac{\pi}{2}<|t|<\pi$ で $-1$ をくり返す方形波（周期 $2\pi$、図\ref{bk:9-2:fig-genten}の右）のフーリエ係数を求めましょう。\par
\kangaekata \kixcbaai{偶関数で、区間ごとに一定}{9-2}{D,B}\par
\kixcdan{1}\enspace $a_n=\frac{1}{\pi}\int_{-\pi}^{\pi}f\cos nt\,\dd t$、$b_n=\frac{1}{\pi}\int_{-\pi}^{\pi}f\sin nt\,\dd t$。\par
\kixcdan{3}\enspace たて軸で折り返すと重なる偶関数なので、$b_n=0$。\par
\kixca $a_n=\frac{2}{\pi}\int_0^{\pi}f\cos nt\,\dd t$（半分の2倍）\par
\kixcnaosu{1}\enspace 0から $\frac{\pi}{2}$ で1、$\frac{\pi}{2}$ から $\pi$ で $-1$ に分ける。\par
\shiki\kixcs $\int_0^{\pi/2}\cos nt\,\dd t=\Bigl[\frac{\sin nt}{n}\Bigr]_0^{\pi/2}=\frac{\sin\frac{n\pi}{2}}{n}$\par
\kixca $\int_{\pi/2}^{\pi}\cos nt\,\dd t=\Bigl[\frac{\sin nt}{n}\Bigr]_{\pi/2}^{\pi}=\frac{0-\sin\frac{n\pi}{2}}{n}$\quad （$\sin n\pi=0$）\par
\kixca $a_n=\frac{2}{\pi}\Bigl(\frac{\sin\frac{n\pi}{2}}{n}+\frac{\sin\frac{n\pi}{2}}{n}\Bigr)=\frac{4\sin\frac{n\pi}{2}}{n\pi}$\quad （2つめは $-1$ 倍）\par
\kixca $\sin\frac{n\pi}{2}$ は $n=1,\ 2,\ 3,\ 4,\ 5$ で $1,\ 0,\ -1,\ 0,\ 1$\par
\kixca $a_0$：上と下の面積が同じなので0（平均0）\par
\kotae $b_n=0$、$a_0=0$、$a_1=\frac{4}{\pi}$、$a_3=-\frac{4}{3\pi}$、$a_5=\frac{4}{5\pi}$、偶数番は0。\\
\hspace*{2.6\zw}$f=\frac{4}{\pi}\bigl(\cos t-\frac13\cos3t+\frac15\cos5t-\cdots\bigr)$
\end{reidai}

\begin{jibun}
$-\frac{\pi}{2}<t<\frac{\pi}{2}$ で1、$\frac{\pi}{2}<|t|<\pi$ で0 をくり返すパルス（周期 $2\pi$）のフーリエ係数を求めましょう。
\bkjans{偶関数なので $b_n=0$。$a_0=\frac{2}{\pi}\int_0^{\pi/2}1\,\dd t=1$ で $\frac{a_0}{2}=\frac12$。$a_n=\frac{2}{\pi}\int_0^{\pi/2}\cos nt\,\dd t=\frac{2\sin\frac{n\pi}{2}}{n\pi}$：$a_1=\frac{2}{\pi}$、$a_2=0$、$a_3=-\frac{2}{3\pi}$、$a_5=\frac{2}{5\pi}$。}
\end{jibun}

\begin{reidai}[山の形の三角波]
$-\pi<t<\pi$ で $f(t)=\pi-|t|$ をくり返す三角波（$t=0$ で高さ $\pi$ の山）のフーリエ係数を求めましょう。\par
\kangaekata \kixcbaai{偶関数で、区間ごとに1次式}{9-2}{D,C}\par
\kixcdan{1}\enspace $a_n=\frac{1}{\pi}\int_{-\pi}^{\pi}(\pi-|t|)\cos nt\,\dd t$ など。\par
\kixcdan{3}\enspace 偶関数なので $b_n=0$、$a_n=\frac{2}{\pi}\int_0^{\pi}(\pi-t)\cos nt\,\dd t$。平均は面積÷$2\pi$。\par
\kixcnaosu{1}\enspace $\pi\cos nt$ の積分と、$t\cos nt$ の部分積分（⑨の三角波と同じ）に分ける。\par
\shiki\kixcs $\frac{a_0}{2}$：山の面積 $\frac12\times2\pi\times\pi=\pi^2$、$\pi^2\div2\pi=\frac{\pi}{2}$\par
\kixca $\int_0^{\pi}\pi\cos nt\,\dd t=\pi\Bigl[\frac{\sin nt}{n}\Bigr]_0^{\pi}=0$\par
\kixca $\int_0^{\pi}t\cos nt\,\dd t=\frac{(-1)^n-1}{n^2}$\quad （⑨で計算した値）\par
\kixca $a_n=\frac{2}{\pi}\Bigl(0-\frac{(-1)^n-1}{n^2}\Bigr)=\frac{2\bigl(1-(-1)^n\bigr)}{\pi n^2}$\par
\kotae $\frac{a_0}{2}=\frac{\pi}{2}$、$b_n=0$、$a_n=\frac{4}{\pi n^2}$（奇数）・0（偶数）。\\
\hspace*{2.6\zw}$f=\frac{\pi}{2}+\frac{4}{\pi}\bigl(\cos t+\frac{\cos3t}{9}+\frac{\cos5t}{25}+\cdots\bigr)$
\end{reidai}

\begin{jibun}
$-\pi<t<\pi$ で $f(t)=2|t|$ をくり返す三角波のフーリエ係数を求めましょう（$\int_0^{\pi}t\cos nt\,\dd t=\frac{(-1)^n-1}{n^2}$ を使ってよい）。
\bkjans{偶関数なので $b_n=0$。面積 $2\times\frac12\times\pi\times2\pi=2\pi^2$、平均 $\frac{a_0}{2}=\pi$。$a_n=\frac{2}{\pi}\int_0^{\pi}2t\cos nt\,\dd t=\frac{4\bigl((-1)^n-1\bigr)}{\pi n^2}$ で、奇数なら $-\frac{8}{\pi n^2}$、偶数なら0。}
\end{jibun}

\begin{reidai}[直線の形の初めの温度]
9-1 の棒（$L=\qty{0.1}{m}$、両はし0\,℃）で、はじめの温度が $f(x)=100\frac{x}{L}$（左はし0\,℃から右はし100\,℃まで直線）のときの、半区間の展開の係数 $b_n$ を求めましょう（右はしは、$t=0$ で急に0\,℃にする）。\par
\kangaekata \kixcbaai{0から $L$ だけの形で、1次式}{9-2}{E,C}\par
\kixcdan{1}\enspace $f=\sum b_n\sin\frac{n\pi x}{L}$ の係数 $b_n$。\par
\kixcdan{3}\enspace 半回転して広げる：$b_n=\frac{2}{L}\int_0^L100\frac{x}{L}\sin\frac{n\pi x}{L}\,\dd x$\par
\kixca $=\frac{200}{L^2}\int_0^Lx\sin\frac{n\pi x}{L}\,\dd x$\quad （定数を外に出す）\par
\kixcnaosu{1}\enspace $x\sin kx$（$k=\frac{n\pi}{L}$）の部分積分。\par
\shiki\kixcs $\int_0^Lx\sin kx\,\dd x=\Bigl[-\frac{x\cos kx}{k}\Bigr]_0^L+\frac{1}{k}\int_0^L\cos kx\,\dd x$\par
\kixca $=-\frac{L\cos n\pi}{k}+\frac{1}{k}\Bigl[\frac{\sin kx}{k}\Bigr]_0^L=-\frac{L(-1)^n}{k}+0$\quad （$kL=n\pi$）\par
\kixca $=\frac{L(-1)^{n+1}}{k}=\frac{L^2(-1)^{n+1}}{n\pi}$\quad （$\frac{1}{k}=\frac{L}{n\pi}$）\par
\kixca $b_n=\frac{200}{L^2}\times\frac{L^2(-1)^{n+1}}{n\pi}=\frac{200(-1)^{n+1}}{n\pi}$\par
\kotae $b_n=\frac{200(-1)^{n+1}}{n\pi}$（$b_1\approx63.66$、$b_2\approx-31.83$、$b_3\approx21.22$）
\end{reidai}

\begin{jibun}
同じ棒で、はじめの温度が $f(x)=50\frac{x}{L}$ のときの $b_n$ を求め、$b_1$、$b_2$、$b_3$ の値を書きましょう。
\bkjans{$b_n=\frac{2}{L}\int_0^L50\frac{x}{L}\sin\frac{n\pi x}{L}\dd x=\frac{100}{L^2}\times\frac{L^2(-1)^{n+1}}{n\pi}=\frac{100(-1)^{n+1}}{n\pi}$。$b_1\approx31.83$、$b_2\approx-15.92$、$b_3\approx10.61$。}
\end{jibun}

\begin{reidai}[まん中ではじいた弦]
長さ $L$ の弦のまん中を高さ $h$ だけつまみ、三角形の形からそっとはなします。はじめの形 $f(x)$（$0\leqq x\leqq\frac{L}{2}$ で $\frac{2h}{L}x$、$\frac{L}{2}\leqq x\leqq L$ で $\frac{2h}{L}(L-x)$）の $b_n$ を求め、どの倍音が出ないかを答えましょう。\par
\kangaekata \kixcbaai{0から $L$ だけの形で、区間ごとに1次式}{9-2}{E,C}\par
\kixcdan{1}\enspace $b_n=\frac{2}{L}\int_0^Lf(x)\sin\frac{n\pi x}{L}\,\dd x$（半区間の展開）。\par
\kixcdan{3}\enspace 右半分は、左半分を $x=\frac{L}{2}$ で折り返した形。$\sin\frac{n\pi x}{L}$ も、$n$ が奇数なら折り返して重なり、偶数なら半回転して重なる。だから右半分の積分は、奇数なら左半分と同じ、偶数なら符号が逆。\par
\kixcnaosu{1}\enspace 左半分 $J=\int_0^{L/2}x\sin kx\,\dd x$（$k=\frac{n\pi}{L}$）を部分積分。\par
\shiki\kixcs $n$ が偶数：左と右が打ち消して $b_n=0$\par
\kixca $n$ が奇数：$b_n=\frac{2}{L}\times\frac{2h}{L}\times2J=\frac{8h}{L^2}J$\par
\kixca $J=\Bigl[-\frac{x\cos kx}{k}\Bigr]_0^{L/2}+\Bigl[\frac{\sin kx}{k^2}\Bigr]_0^{L/2}=-\frac{L\cos\frac{n\pi}{2}}{2k}+\frac{\sin\frac{n\pi}{2}}{k^2}$\par
\kixca 奇数では $\cos\frac{n\pi}{2}=0$ なので、$J=\frac{\sin\frac{n\pi}{2}}{k^2}=\frac{L^2\sin\frac{n\pi}{2}}{n^2\pi^2}$\par
\kixca $b_n=\frac{8h}{L^2}\times\frac{L^2\sin\frac{n\pi}{2}}{n^2\pi^2}=\frac{8h\sin\frac{n\pi}{2}}{n^2\pi^2}$\par
\kotae $b_n=\frac{8h}{n^2\pi^2}\sin\frac{n\pi}{2}$（$b_1\approx0.811h$、$b_3\approx-0.090h$、$b_5\approx0.032h$）。偶数番の倍音は出ない
\end{reidai}

\begin{jibun}
9-1 の棒で、はじめの温度が三角形（両はし0\,℃、まん中100\,℃、まん中まで直線）のときの $b_n$ を求め、$b_1$、$b_2$、$b_3$ の値を書きましょう（例題5の $h$ を100にした形）。
\bkjans{例題5と同じ計算で $b_n=\frac{800}{n^2\pi^2}\sin\frac{n\pi}{2}$。$b_1\approx81.06$、$b_2=0$、$b_3\approx-9.006$。}
\end{jibun}

\begin{reidai}[50\,Hz の方形波を $\frac{2}{T}$ の式で]
周期 $T=\qty{0.02}{s}$（50\,Hz）で、$0<t<\frac{T}{2}$ で1、$\frac{T}{2}<t<T$ で $-1$ をくり返す方形波の係数を、$\frac{2}{T}\int$ の式で求め、成分の周波数を答えましょう。\par
\kangaekata \kixcbaai{周期が $2\pi$ でない、奇関数、区間ごとに一定}{9-2}{F,D,B}\par
\kixcdan{1}\enspace $b_n=\frac{2}{T}\int_0^Tf\sin n\omega_0t\,\dd t$、$\omega_0=\frac{2\pi}{T}=100\pi$\,rad/s。\par
\kixcdan{3}\enspace 原点で半回転すると重なる奇関数なので、$a_n=0$。\par
\kixca $b_n=\frac{4}{T}\int_0^{T/2}1\times\sin n\omega_0t\,\dd t$（半周期の2倍）\par
\kixcnaosu{1}\enspace $\sin n\omega_0t$ の定積分。\par
\shiki\kixcs $\int_0^{T/2}\sin n\omega_0t\,\dd t=\Bigl[-\frac{\cos n\omega_0t}{n\omega_0}\Bigr]_0^{T/2}=\frac{1-\cos\frac{n\omega_0T}{2}}{n\omega_0}$\par
\kixca $\omega_0T=2\pi$ なので $\frac{n\omega_0T}{2}=n\pi$、分子は $1-(-1)^n$\par
\kixca $b_n=\frac{4}{T}\times\frac{1-(-1)^n}{n\omega_0}=\frac{4\bigl(1-(-1)^n\bigr)}{n\times2\pi}=\frac{2\bigl(1-(-1)^n\bigr)}{n\pi}$\quad （$T\omega_0=2\pi$）\par
\kixca 奇数なら $\frac{4}{n\pi}$、偶数なら0。$n$ 番めの周波数は $50n$\,Hz\par
\kotae $b_n=\frac{4}{n\pi}$（奇数）。50\,Hz：1.27、150\,Hz：0.42、250\,Hz：0.25（9-1④と同じ）
\end{reidai}

\begin{jibun}
周期 $T=\frac{1}{60}$\,s（60\,Hz）で、前半が5\,V、後半が $-5$\,V をくり返す電圧の係数を、$\frac{2}{T}\int$ の式で求め、60\,Hz・180\,Hz・300\,Hz の成分の大きさを答えましょう。
\bkjans{奇関数なので $a_n=0$。$b_n=\frac{4}{T}\int_0^{T/2}5\sin n\omega_0t\,\dd t=\frac{10\bigl(1-(-1)^n\bigr)}{n\pi}$ で、奇数なら $\frac{20}{n\pi}$。60\,Hz：約6.37\,V、180\,Hz：約2.12\,V、300\,Hz：約1.27\,V。}
\end{jibun}

\begin{reidai}[整流した波の平均]
半波整流した波（$0<t<\pi$ で $\sin t$、$\pi<t<2\pi$ で0、周期 $2\pi$）の平均を求めましょう。また、もとの交流の最大値が \qty{141}{V}（家の100\,V の交流の山の高さ）のとき、平均の電圧は何Vですか。\par
\kangaekata \kixcbaai{平均（直流分）だけほしい}{9-2}{G}\par
\kixcdan{1}\enspace $\frac{a_0}{2}=\frac{1}{2\pi}\int_0^{2\pi}f\,\dd t$。\par
\kixcdan{3}\enspace 1周期の面積÷周期。面積は山1つぶん。\par
\shiki\kixcs 山1つの面積 $\int_0^{\pi}\sin t\,\dd t=\Bigl[-\cos t\Bigr]_0^{\pi}=1-(-1)=2$\quad （5-15）\par
\kixca 平均 $2\div2\pi=\frac{1}{\pi}\approx0.318$\par
\kixca $141\times\frac{1}{\pi}\approx44.9$\par
\begin{center}\begin{tikzpicture}
\begin{axis}[bkplot, width=100mm, height=32mm, axis x line=bottom, axis y line=left, xmin=0, xmax=740, ymin=0, ymax=1.2, xtick={180,360,540,720}, xticklabels={$\pi$,$2\pi$,$3\pi$,$4\pi$}, ytick={0.318,1}, yticklabels={$\frac{1}{\pi}$,1}, xlabel={$t$}]
  \addplot[bkBlue, line width=1.2pt, domain=0:720, samples=300] {max(sin(x),0)};
  \addplot[bkOrange!85!black, line width=0.9pt, dash pattern=on 3pt off 1.5pt] coordinates {(0,0.318) (720,0.318)};
\end{axis}
\end{tikzpicture}\end{center}
\kixca 青が半波整流した波、オレンジの点線が平均。山の面積を2周期ぶんにならした高さ\par
\kotae 平均 $\frac{1}{\pi}\approx0.318$（最大値の約32\%）。141\,V なら約44.9\,V
\end{reidai}

\begin{jibun}
全波整流した波（$|\sin t|$。周期 $\pi$ で山がならぶ）の平均を求めましょう。最大値が \qty{141}{V} なら、平均の電圧は何Vですか。
\bkjans{山1つの面積は2、周期は $\pi$ なので、平均 $\frac{2}{\pi}\approx0.637$。$141\times\frac{2}{\pi}\approx89.8$\,V。}
\end{jibun}

\begin{reidai}[半波整流の $\cos2t$ の係数]
半波整流した波（例題7）の、$a_2$ と $b_1$ を求めましょう。\par
\kangaekata \kixcbaai{区間ごとに $\sin$ の形}{9-2}{H}\par
\kixcdan{1}\enspace $a_2=\frac{1}{\pi}\int_0^{2\pi}f\cos2t\,\dd t$、$b_1=\frac{1}{\pi}\int_0^{2\pi}f\sin t\,\dd t$。\par
\kixcnaosu{1}\enspace $\pi<t<2\pi$ では $f=0$ なので、0から $\pi$ だけ。積和の式で和に直す。\par
\shiki\kixcs $\sin t\cos2t=\frac12\bigl\{\sin3t+\sin(-t)\bigr\}=\frac12(\sin3t-\sin t)$\par
\kixca $\int_0^{\pi}\sin3t\,\dd t=\frac{1-\cos3\pi}{3}=\frac23$、\quad $\int_0^{\pi}\sin t\,\dd t=2$\par
\kixca $a_2=\frac{1}{\pi}\times\frac12\Bigl(\frac23-2\Bigr)=\frac{1}{2\pi}\times\Bigl(-\frac43\Bigr)=-\frac{2}{3\pi}$\par
\kixca $b_1=\frac{1}{\pi}\int_0^{\pi}\sin^2t\,\dd t=\frac{1}{\pi}\times\frac{\pi}{2}=\frac12$\quad （$\sin^2$ の平均 $\frac12$×長さ $\pi$）\par
\kixca 同じように計算した係数を並べると、下のスペクトルになる（$a_4$ はじぶんでとこう）\par
\begin{center}\begin{tikzpicture}[x=1mm, y=1mm]
  \draw[bkInk, line width=0.5pt, -{Stealth[length=1.6mm]}] (0,0) -- (78,0) node[right, font=\footnotesize, inner sep=1pt] {番号};
  \foreach \kixcx/\kixch/\kixcl in {6/0.3183/{平均},20/0.5/{$b_1$},34/0.2122/{$|a_2|$},48/0.0424/{$|a_4|$},62/0.0182/{$|a_6|$}} {%
    \fill[bkBlue!80] ({\kixcx-2},0) rectangle ({\kixcx+2},{30*\kixch});
    \node[font=\footnotesize, anchor=north] at (\kixcx,-0.8) {\kixcl};}
  \node[font=\footnotesize, anchor=south] at (6,9.8) {0.318};
  \node[font=\footnotesize, anchor=south] at (20,15.2) {0.5};
  \node[font=\footnotesize, anchor=south] at (34,6.6) {0.212};
\end{tikzpicture}\end{center}
\kotae $a_2=-\frac{2}{3\pi}\approx-0.212$、$b_1=\frac12$
\end{reidai}

\begin{jibun}
半波整流した波の $a_4$ を求めましょう。
\bkjans{$\sin t\cos4t=\frac12(\sin5t-\sin3t)$。$\int_0^{\pi}\sin5t\,\dd t=\frac25$、$\int_0^{\pi}\sin3t\,\dd t=\frac23$。$a_4=\frac{1}{\pi}\times\frac12\bigl(\frac25-\frac23\bigr)=-\frac{2}{15\pi}\approx-0.0424$。}
\end{jibun}

\begin{reidai}[12個のデータから係数を求める]
ある機械の振動を、1周期の中で30° ごとに12回はかりました（$t$ は角で表した時刻。$f$ の単位は mm）。$\sin t$ の係数 $b_1$ を求めましょう。\par
\begin{center}{\footnotesize\begin{tblr}{colspec={Q[l,m,9mm] *{12}{Q[c,m,7mm]}}, hlines={0.4pt,bkLine}, vlines={0.4pt,bkLine}, rowsep=0.5pt, colsep=1pt, row{1}={bg=bkLine!40}}
$t$（度） & 0 & 30 & 60 & 90 & 120 & 150 & 180 & 210 & 240 & 270 & 300 & 330 \\
$f$ & 1.0 & 2.0 & 2.1 & 2.0 & 2.1 & 2.0 & 1.0 & $-1.0$ & $-3.1$ & $-4.0$ & $-3.1$ & $-1.0$ \\
\end{tblr}}\end{center}
\kangaekata \kixcbaai{式がなく、はかった値だけ}{9-2}{I}\par
\kixcdan{1}\enspace $b_1=\frac{1}{\pi}\int_0^{2\pi}f\sin t\,\dd t$。式がないので積分できない。\par
\kixcdan{3}\enspace $b_1\approx\frac{2}{12}\sum f_k\sin t_k$（12個の「値×$\sin$」の平均の2倍）。\par
\shiki\kixcs $\sin t_k$：0、0.5、0.866、1、0.866、0.5、0、$-0.5$、$-0.866$、$-1$、$-0.866$、$-0.5$\par
\kixca 積 $f_k\sin t_k$：0、1.0、1.819、2.0、1.819、1.0、0、0.5、2.685、4.0、2.685、0.5\par
\kixca 合計 18.008\par
\kixca $b_1\approx\frac{2}{12}\times18.008=3.001$\par
\kotae $b_1\approx\qty{3.00}{mm}$
\end{reidai}

\begin{jibun}
例題9のデータから、$\cos2t$ の係数 $a_2$ を求めましょう（$\cos2t_k$ は、1、0.5、$-0.5$、$-1$、$-0.5$、0.5、1、0.5、$-0.5$、$-1$、$-0.5$、0.5）。
\bkjans{積 $f_k\cos2t_k$：1.0、1.0、$-1.05$、$-2.0$、$-1.05$、1.0、1.0、$-0.5$、1.55、4.0、1.55、$-0.5$。合計6.0。$a_2\approx\frac{2}{12}\times6.0=1.00$（mm）。}
\end{jibun}

\begin{reidai}[方形波のひずみ率]
高さ1の方形波（$b_n=\frac{4}{n\pi}$、奇数番）のひずみ率を、パーセバルの等式を使って求めましょう。\par
\kangaekata 係数はわかっているので、表9-2 は使わない。パーセバルの等式で、高調波を無限個たさずにすませる。\par
\kixcdan{1}\enspace ひずみ率$=\sqrt{\frac{\text{高調波の電力の合計}}{\text{基本波の電力}}}$。電力は「係数の2乗の $\frac12$」。\par
\shiki\kixcs 全体の電力＝2乗の平均＝1（方形波の2乗はいつも1）\par
\kixca 基本波の電力 $\frac12\Bigl(\frac{4}{\pi}\Bigr)^2=\frac{8}{\pi^2}\approx0.8106$\par
\kixca 高調波の電力の合計 $1-0.8106=0.1894$\quad （パーセバル：全体は成分の和）\par
\kixca $\frac{0.1894}{0.8106}=0.2337$、$\sqrt{0.2337}\approx0.483$\par
\kotae 約48.3\%
\end{reidai}

\begin{jibun}
三角波 $|t|$（$-\pi<t<\pi$。$a_1=-\frac{4}{\pi}$）のひずみ率を、パーセバルの等式で求めましょう。平均 $\frac{\pi}{2}$ をひいた $|t|-\frac{\pi}{2}$ の2乗の平均は $\frac{\pi^2}{12}\approx0.8225$ です。
\bkjans{基本波の電力 $\frac12\bigl(\frac{4}{\pi}\bigr)^2\approx0.8106$。高調波の電力 $0.8225-0.8106=0.0119$。$\sqrt{\frac{0.0119}{0.8106}}=\sqrt{0.0147}\approx0.121$。約12.1\%。}
\end{jibun}

%% ---------------- まちがえやすいところ ----------------
\bksection{tsumazuki}{まちがえやすいところをチェック}

\begin{tsumazuki}{① 半周期にしたのに、2倍し忘れる}
\bkng{偶関数の $a_n=\frac{1}{\pi}\int_0^{\pi}f\cos nt\,\dd t$}
\bkok{$a_n=\frac{2}{\pi}\int_0^{\pi}f\cos nt\,\dd t$（半分の区間の2倍）}
\[ \kixcNt{a_n}=\kixcSt{1}{\frac{1}{\pi}\int_{-\pi}^{\pi}f\cos nt\,\dd t}=\kixcSt{1}{\frac{2}{\pi}\int_0^{\pi}f\cos nt\,\dd t} \]
D で楽をするのは「区間を半分にして2倍する」ことです。名前の式（1周期）から、1行ずつ直せばまちがえません。
\end{tsumazuki}

\begingroup\tcbset{equal height group=kixcEHA}%
\noindent\begin{minipage}[t]{0.485\linewidth}\vspace{0pt}
\begin{tsumazuki}{② 定数の項を $a_0$ と書く}
\raggedright\bkng{例題3の定数の項は $a_0=\pi$}
\bkok{$a_0=\pi$ だが、定数の項は $\frac{a_0}{2}=\frac{\pi}{2}$}
1の「長さの2乗」は $T$ で、ほかの成分の2倍です。だから定数の項だけ $\frac12$ がつきます。
\end{tsumazuki}
\end{minipage}\hfill
\begin{minipage}[t]{0.485\linewidth}\vspace{0pt}
\begin{tsumazuki}{③ 偶関数・奇関数を取りちがえる}
\raggedright\bkng{$|t|$ は奇関数、0と1の信号は奇関数}
\bkok{$|t|$ は偶関数。0と1の信号はどちらでもない}
$f(-t)$ を計算して確かめます。方形波は、原点の置き方で変わります。
\end{tsumazuki}
\end{minipage}\par
\endgroup

\begingroup\tcbset{equal height group=kixcEHB}%
\noindent\begin{minipage}[t]{0.485\linewidth}\vspace{0pt}
\begin{tsumazuki}{④ 半区間の展開で $\cos$ を使う}
\raggedright\bkng{棒の温度を $\cos\frac{n\pi x}{L}$ でたす}
\bkok{両はし0なら $\sin\frac{n\pi x}{L}$ だけ}
原点で半回転して広げた形は奇関数なので、$\cos$ の係数はすべて0です。$\cos$ は両はしで0になりません。
\end{tsumazuki}
\end{minipage}\hfill
\begin{minipage}[t]{0.485\linewidth}\vspace{0pt}
\begin{tsumazuki}{⑤ $\sin\frac{n\pi}{2}$ を $\sin n\pi$ と同じにする}
\raggedright\bkng{$\sin\frac{n\pi}{2}=0$}
\bkok{$\sin\frac{n\pi}{2}$ は $1,\ 0,\ -1,\ 0,\ \cdots$}
$\sin n\pi$ はいつも0ですが、$\sin\frac{n\pi}{2}$ は奇数番で $\pm1$ です。
\end{tsumazuki}
\end{minipage}\par
\endgroup

\begingroup\tcbset{equal height group=kixcEHC}%
\noindent\begin{minipage}[t]{0.485\linewidth}\vspace{0pt}
\begin{tsumazuki}{⑥ データの和で $\frac{2}{N}$ を $\frac{1}{N}$ にする}
\raggedright\bkng{$b_1\approx\frac{1}{12}\sum f_k\sin t_k$}
\bkok{$b_1\approx\frac{2}{12}\sum f_k\sin t_k$}
係数は「平均の2倍」です。平均（$\frac{1}{N}\sum$）だけでは半分の値になります。
\end{tsumazuki}
\end{minipage}\hfill
\begin{minipage}[t]{0.485\linewidth}\vspace{0pt}
\begin{tsumazuki}{⑦ 成分の電力を係数の2乗にする}
\raggedright\bkng{基本波の電力 $b_1^2$}
\bkok{基本波の電力 $\frac12b_1^2$}
$\sin$ の2乗の平均は $\frac12$ です。パーセバルの等式にも $\frac12$ がつきます。
\end{tsumazuki}
\end{minipage}\par
\endgroup

%% ---------------- 物理・工学ではこう使う ----------------
\begin{butsuri}
\kixcbox{弱い信号を取り出す：ロックインアンプ}
\noindent\begin{minipage}[c]{0.64\linewidth}
\raggedright
はてな？の「ほしい周波数の $\sin$ をかけて平均する」は、そのまま計測器になっています。ロックインアンプは、はかりたい信号に、決まった周波数の $\sin$ を電子回路でかけて平均し、その周波数の成分だけを取り出します。ほかの周波数の雑音は、直交性で平均が0になって消えるので、雑音に100倍うもれた小さな信号でも取り出せることがあります。
\end{minipage}\hfill
\begin{minipage}[c]{0.32\linewidth}
\centering
\begin{tikzpicture}[x=1mm, y=1mm]
  \node[kixc do, text width=12mm] (a) at (0,0) {信号\\＋雑音};
  \node[kixc do, text width=10mm] (b) at (20,0) {×$\sin$};
  \node[kixc do, text width=10mm] (c) at (20,-14) {平均};
  \draw[kixc go] (a) -- (b);
  \draw[kixc go] (b) -- (c);
  \node[font=\footnotesize, anchor=north] at (14,-19) {ほしい成分だけ};
\end{tikzpicture}
\end{minipage}
\par\vspace{6pt}
\kixcbox{電気のひずみをはかる}
工場やビルの電気設備では、電流の波形をはかる計器（電力計・ひずみ率計）が、1周期のデータからフーリエ係数を計算し、第3・第5・第7高調波の大きさとひずみ率を表示します（場合 I とパーセバルの等式）。整流をする機器（パソコンの電源、インバーター）は、高調波の多い電流を流すので、高調波を小さくするフィルターをつけます。
\par\vspace{6pt}
\kixcbox{音のイコライザーと、機械の振動}
音楽アプリのイコライザーは、音を周波数ごとの成分に分け、低音・高音の成分の大きさを変えてから、たし直します。回る機械の振動を周波数ごとに分けると、回転の何倍の成分が大きいかで、軸のずれや羽根のこわれなど、どこが悪いかの見当がつきます（9-4）。
\end{butsuri}

\begin{column}{イコライザー：成分ごとに大きさを変えて、たし直す}
\noindent\begin{minipage}[c]{0.56\linewidth}
\raggedright
音楽アプリや音響機器のイコライザーは、音を周波数の帯（たとえば 31\,Hz・63\,Hz・125\,Hz…16\,kHz の10の帯。となりの帯は周波数が2倍）に分け、帯ごとに大きさを何倍かしてから、たし直します。フーリエ級数でいえば、係数 $a_n$・$b_n$ に帯ごとの倍率をかけて、もう一度 $\sum$ でたすことです。\par\vspace{3pt}
低い音を2倍にすると、2乗の平均（パーセバルの等式）で、その帯の電力は4倍になります。だから、低音を上げすぎると、スピーカーやアンプに大きな電力がかかります。
\end{minipage}\hfill
\begin{minipage}[c]{0.4\linewidth}
\centering
\begin{tikzpicture}[x=0.82mm, y=1mm]
  \draw[bkInk, line width=0.5pt, -{Stealth[length=1.6mm]}] (0,0) -- (52,0) node[right, font=\footnotesize, inner sep=1pt] {帯};
  \foreach \kixcx/\kixca/\kixcb in {4/6/12,12/8/12,20/7/7,28/6/6,36/5/5,44/4/6} {%
    \fill[kixcF!35] ({\kixcx-2.6},0) rectangle ({\kixcx-0.2},\kixca);
    \fill[kixcW] ({\kixcx+0.2},0) rectangle ({\kixcx+2.6},\kixcb);}
  \fill[kixcF!35] (2,18) rectangle (5,20.5); \node[font=\footnotesize, anchor=west] at (5.5,19.25) {もと};
  \fill[kixcW] (24,18) rectangle (27,20.5); \node[font=\footnotesize, anchor=west] at (27.5,19.25) {変えたあと};
  \node[font=\footnotesize, anchor=north] at (8,-1) {低い音};
  \node[font=\footnotesize, anchor=north] at (44,-1) {高い音};
\end{tikzpicture}\par
{\footnotesize 低い2つの帯を2倍、いちばん高い帯を1.5倍にした例}
\end{minipage}
\end{column}


%% ---------------- 章末問題 ----------------
\bkendproblems
\bkproblemgroup{基本}{5}

\bkproblem $f(t)=\sin t\cos^2t$ のフーリエ係数を求めましょう（周期 $2\pi$）\mbox{（\bkreidairef{1}）}。
\bkpans{$\cos^2t=\frac{1+\cos2t}{2}$ から $f=\frac12\sin t+\frac12\sin t\cos2t$。$\sin t\cos2t=\frac12(\sin3t-\sin t)$ なので、$f=\frac14\sin t+\frac14\sin3t$。$b_1=\frac14$、$b_3=\frac14$、ほかは0。}

\bkproblem $-\frac{\pi}{2}<t<\frac{\pi}{2}$ で3、$\frac{\pi}{2}<|t|<\pi$ で0 をくり返すパルス（周期 $2\pi$）のフーリエ係数を求めましょう\mbox{（\bkreidairef{2}）}。
\bkpans{偶関数なので $b_n=0$。$\frac{a_0}{2}=\frac32$。$a_n=\frac{2}{\pi}\int_0^{\pi/2}3\cos nt\,\dd t=\frac{6\sin\frac{n\pi}{2}}{n\pi}$：$a_1=\frac{6}{\pi}$、$a_2=0$、$a_3=-\frac{2}{\pi}$。}

\bkproblem $-\pi<t<\pi$ で $f(t)=\frac{\pi}{2}-|t|$ をくり返す三角波のフーリエ係数を求めましょう\mbox{（\bkreidairef{3}）}。
\bkpans{偶関数で $b_n=0$。平均（面積）は0 で $\frac{a_0}{2}=0$。$a_n=\frac{2}{\pi}\int_0^{\pi}\bigl(\frac{\pi}{2}-t\bigr)\cos nt\,\dd t=\frac{2\bigl(1-(-1)^n\bigr)}{\pi n^2}$ で、奇数なら $\frac{4}{\pi n^2}$、偶数なら0。}

\bkproblem 9-1 の棒（$L=\qty{0.1}{m}$）で、はじめの温度が $f(x)=20\frac{x}{L}$ のときの $b_n$ と、$b_1$、$b_2$ の値を求めましょう\mbox{（\bkreidairef{4}）}。
\bkpans{$b_n=\frac{40(-1)^{n+1}}{n\pi}$。$b_1\approx12.73$、$b_2\approx-6.37$。}

\bkproblem 周期 \qty{1}{ms}（1\,kHz）で、前半が2\,V、後半が $-2$\,V をくり返す電圧の、1\,kHz・3\,kHz・5\,kHz の成分の大きさを、$\frac{2}{T}\int$ の式で求めましょう\mbox{（\bkreidairef{6}）}。
\bkpans{奇関数で $a_n=0$。$b_n=\frac{4}{T}\int_0^{T/2}2\sin n\omega_0t\,\dd t=\frac{4\bigl(1-(-1)^n\bigr)}{n\pi}$、奇数なら $\frac{8}{n\pi}$。1\,kHz：約2.55\,V、3\,kHz：約0.85\,V、5\,kHz：約0.51\,V。}

\bkproblem 最大値が \qty{283}{V}（200\,V の交流）の電圧を半波整流しました。平均の電圧は何Vですか\mbox{（\bkreidairef{7}）}。
\bkpans{平均は最大値の $\frac{1}{\pi}$ 倍。$283\times\frac{1}{\pi}\approx90.1$\,V。}

\bkproblemgroup{チャレンジ}{5}

\bkproblem 半波整流した波の $a_6$ を求めましょう\mbox{（\bkreidairef{8}）}。
\bkpans{$\sin t\cos6t=\frac12(\sin7t-\sin5t)$。$\int_0^{\pi}\sin7t\,\dd t=\frac27$、$\int_0^{\pi}\sin5t\,\dd t=\frac25$。$a_6=\frac{1}{2\pi}\bigl(\frac27-\frac25\bigr)=-\frac{2}{35\pi}\approx-0.0182$。}

\bkproblem 例題9のデータから、$\sin2t$ の係数 $b_2$ を求めましょう（$\sin2t_k$ は、0、0.866、0.866、0、$-0.866$、$-0.866$、0、0.866、0.866、0、$-0.866$、$-0.866$）\mbox{（\bkreidairef{9}）}。
\bkpans{積 $f_k\sin2t_k$：0、1.732、1.819、0、$-1.819$、$-1.732$、0、$-0.866$、$-2.685$、0、2.685、0.866。合計0。$b_2=0$。}

\bkproblem 0と1をくり返す信号（9-1 の例題2。$b_n=\frac{2}{n\pi}$、奇数番）のひずみ率を、パーセバルの等式で求めましょう。平均 $\frac12$ をひいた形の2乗の平均は $\frac14$ です\mbox{（\bkreidairef{10}）}。
\bkpans{基本波の電力 $\frac12\bigl(\frac{2}{\pi}\bigr)^2=\frac{2}{\pi^2}\approx0.2026$。高調波の電力 $0.25-0.2026=0.0474$。$\sqrt{\frac{0.0474}{0.2026}}\approx\sqrt{0.234}\approx0.483$。約48.3\%（方形波と同じ）。}

\begin{furikaeri}
\begin{checklist}
\item ちがう周波数の $\sin$ どうしの積を1周期でたすと0になることを、表計算とグラフで確かめられる。\bkfurref{bk:9-2:kakeru}
\item 直交性を使って、係数が「かけて積分し、長さの2乗でわる」で取り出せるわけを説明できる。\bkfurref{bk:9-2:toridasu}
\item 場合分けの表9-2 と流れ図で、波の形に合った楽な求め方を選び、係数を計算できる。\bkfurref{bk:9-2:basho}
\end{checklist}
\end{furikaeri}

\bkfukidashi{mina}{$\sin$ と $\cos$ を2つ使うと、式が長くなるね。次の 9-3 では、6-10 の $e^{j\theta}$ を使って、2つを1つにまとめてみよう。}
